% !TEX program = xelatex
% ============================================================
%  Philosophy of Intimacy and the Theory of Justice · Paper VII
%  Prolegomena to a Diplomacy of Intimate Relationships:
%  A First Step toward Generalized Generative Relational Being
%  — Value Generation, Justice, and the Gaze of the Other
%    in the External Relational Field
%
%  Series visual identity: blush page, deep-rose headings,
%  rose emphasis, gold rules, ink body. XeLaTeX + TeX Gyre Pagella,
%  with Noto Serif CJK SC fallback for Chinese terms.
% ============================================================
\documentclass[11pt,a4paper]{article}
\PassOptionsToPackage{table}{xcolor}

% ---------- Fonts (CJK fallback for occasional Chinese terms) ----------
\usepackage{fontspec}
\newfontfamily\cjk{Noto Serif CJK SC}
\newcommand{\zh}[1]{{\cjk #1}}
\XeTeXlinebreaklocale "zh"
\XeTeXlinebreakskip = 0pt plus 0.1pt

% natbib (numbered) before serendip-paper (hyperref loads inside it)
\usepackage[numbers,sort&compress]{natbib}
\usepackage{serendip-paper}

% colour used on the title page (headings/links come from serendip-paper)
\definecolor{inksoft}{HTML}{342E37}

% ---------- Tables (beyond the shared style) ----------
\usepackage{tabularx}
\usepackage{pdflscape}
\usepackage{ragged2e}
\renewcommand{\arraystretch}{1.3}
\newcolumntype{Y}{>{\RaggedRight\arraybackslash}p{3.1cm}}

% ---------- Paper-specific emphasis & helpers ----------
\newcommand{\emc}[1]{\textcolor{rose}{\oldtextit{#1}}}
\newcommand{\emb}[1]{\textcolor{rose}{\oldtextbf{#1}}}
\newcommand{\goldrule}{\noindent\textcolor{gold}{\rule{\linewidth}{0.6pt}}\par\vspace{0.4em}}
\newcommand{\todo}[1]{\textcolor{gray}{\small[TODO: #1]}}
\newcommand{\xref}[1]{\textcolor{rose}{\S\ref{#1}}}
\newcommand{\pinkref}[1]{\textcolor{rose}{#1}}
\newenvironment{casebox}[1]{%
  \par\vspace{0.4em}\noindent\textit{\textcolor{rose}{An illustrative case --- #1.}}\hspace{0.4em}\ignorespaces}{%
  \par\vspace{0.4em}}

% ============================================================
\begin{document}

% ---------- Title page ----------
\thispagestyle{empty}
\begin{center}
\vspace*{2cm}
{\footnotesize\color{rosedeep}\scshape Philosophy of Intimacy and the Theory of Justice \quad$\cdot$\quad Paper VII}\\[0.3cm]
\coverrule\\[1.2cm]
{\LARGE\bfseries\color{warnred} Prolegomena to a Diplomacy of Intimate Relationships}\\[0.7cm]
{\large\color{warnred} A First Step toward Generalized Generative Relational Being}\\[0.8cm]
{\small\color{gold} [\,Working Draft\,]}\\[0.8cm]
{\normalsize\color{inksoft} Value Generation, Justice, and the Gaze of the Other\\[0.15cm] in the External Relational Field}\\[2cm]
\coverrule[0.25]\\[0.6cm]
{\normalsize\color{inksoft} Wanhong}\\
\vfill
{\footnotesize\color{gold} Working draft --- not for citation or circulation}
\coverfootnote
\end{center}
\newpage

% ---------- Epigraph ----------
\thispagestyle{empty}
\vspace*{3cm}
\begin{center}
\begin{minipage}{0.72\linewidth}
\centering\itshape\color{rose}
{\cjk 此有故彼有，此生故彼生。}\\[0.5em]
This being, that is; this arising, that arises.\\[0.4em]
{\upshape\footnotesize\color{gold} ---{\cjk 《杂阿含经》}\quad \emph{Sa\d{m}yukta {\=A}gama}}
\end{minipage}
\end{center}
\newpage

\tableofcontents
\newpage

% ============================================================
\section{Introduction: From the Dyad to the World}

The first six papers of this series have, throughout, worked inside a tacitly closed system. The mother paper established the ontology of relational being: that being is set apart, prior to relation, is an illusion of modernity; the truth is that the subject is generated in relation, through relation. \pinkref{Paper~III} treated the normativity of this relational system, arguing that the norms of an intimate bond are not imposed upon the two parties from without but are the product of their joint self-legislation. \pinkref{Paper~V} descended to the micro-mechanism of subject formation, using joint attention to characterize how two consciousnesses mutually constitute one another within a shared attentional field. \pinkref{Paper~VI} cut into a single decisive moment---the proposal of marriage---and exposed the epistemic injustice latent within it: one party may, procedurally, solicit consent in full, yet substantively monopolize the authority to interpret the content of ``love,'' ``commitment,'' and ``we.''

These six papers share a premise that goes unstated: they all unfold at the scale of the \emb{dyad}. Whether the matter is self-legislation, joint attention, or epistemic recognition, the stage of the argument holds only two persons. The external world---the parents, siblings, old friends, and colleagues of either party, and the anonymous social norms that surround them---appears at most as a backdrop, never as a \emb{coupling term} entering the system. This closure was no oversight but a methodological simplification of necessity: first make the mechanisms clear within the smallest structurally complete relational system, then generalize outward.

\emc{The present paper is that first step of generalization outward.} Its central claim is this: \emb{from the moment of its formation, the dyad is already surrounded, watched, and normed by a network of external relations; the external relational field is not an optional addition to the dyad but a constitutive condition of its existence.} The ``we'' never establishes itself in a vacuum; it is given, recognized, and contested within a field already dense with others. When a couple first walks hand in hand through the door of either family's home, they are already engaged in an activity---presenting themselves, negotiating a position, petitioning for acceptance---before an authority that exists prior to them and that holds the power of recognition and the disposal of resources. This activity the present paper calls \emb{diplomacy}.

To call it diplomacy marks three facts. First, in the external relational field the dyad is not an isolated atom but a \emb{quasi-sovereign unit}: it has an inside and an outside, a common position that must be represented abroad, internal affairs that must be guarded. Second, its dealings with external relational terms are \emb{structured, regulated, and freighted with stakes}---arrivals call for rites, exchanges incur debts, breaches of decorum carry consequences, much as between states. Third, and most decisively, this commerce is \emb{often not diplomacy between equals but asymmetric diplomacy}: before parents and lineage, the dyad is frequently not an equal contracting party but rather the lower-standing side presenting its credentials to an authority that holds history, resources, and the power of definition (see \xref{sec:power}).

In the vocabulary of this series, the leap is from \emb{second-order coupling} to \emb{higher-order coupling}. Second-order coupling is the mutual constitution of two subjects; higher-order coupling is the mutual constitution between ``the dyad as a unit'' and one or more external relational terms. This is not merely raising the head-count from two to three, four, five. Higher-order coupling brings \emb{structurally new phenomena} that simply do not exist at the dyadic scale: the formation of triangles, the objectification of being-watched-as-a-whole, the ascent into nested mentalizing, the circulation and extraction of value within a larger network. These will be brought to light, one by one, below. The series names this larger research programme \emb{Generalized Generative Relational Being} (GRB): generative, because relation does not rest in static connection but continually produces value, subjects, and norms; generalized, because it must extend from the dyad to relational structures of arbitrary order. The present paper is the first step toward that programme.

\emc{The standing of this paper must be clarified at the outset.} This is a \emb{programmatic review}, not a finished argument. Scholarly writing ordinarily expects a paper to advance one thesis, mount one line of argument, reach one conclusion. This is not such a paper. Its task is \emb{mapping}: within the vast and almost wholly unsurveyed problem-domain of ``diplomacy,'' to mark out which sub-problems there are, to which theoretical stratum each belongs, how they couple with one another, and to erect a coordinate system and a literature infrastructure for the sub-series of papers that will each go deeper in turn. It seeks breadth, not exhaustion. Wherever a matter could be argued at depth, this paper gives the minimal analytic frame and then marks ``to be developed in a later paper,'' indicating its destination in the roadmap of the conclusion (\xref{sec:roadmap}).

This standing dictates two features of the paper's apparatus. First, it makes heavy use of \emb{literature tables}: each problem-domain is not only analyzed but accompanied by its key sources, their core contributions, and---crucially---their specific significance for this research programme, so that the paper serves as an entry map for later researchers. Second, beyond the abstract coupling analysis, it embeds \emb{thick case descriptions}: at several nodes of the relational network it places a concrete, anonymized situation, so that a phenomenon shows itself in its particularity rather than remaining a bare name on a coordinate axis. The methodological account of these two devices is given in \xref{sec:method}.

Finally, a word is owed on why this paper does \emb{not} adopt a seemingly more natural manner of writing---taking some one phenomenon (the most conspicuous candidate being the gift) as a hub running through the whole. This refusal is no matter of stylistic preference but a methodological requirement derived from within the relational ontology itself. The next section takes up the matter.

% ============================================================
\section{A Methodological Self-Examination: Why No Central Phenomenon}

Consider a tempting manner of writing. Among all the phenomena of diplomacy, the gift seems to sit naturally at the center: it is weighty, visible, measurable, and nearly every abstract dimension shows itself at once upon it---it has value (the theory of value), it is a sign (semiotics), it must be read (hermeneutics), it manufactures indebtedness (the theory of justice), it transfers wealth across kin networks (political economy), and it is always given and received under the gaze of others (the gaze). Since the gift is so richly endowed, why not take it as the hub, let the whole paper unfold around it, and treat the remaining phenomena as its footnotes or extensions?

This paper \emb{refuses} that manner. There are three reasons, and all three grow from within the relational ontology of this series---which makes the refusal itself a theoretical stance rather than a matter of taste.

\emc{First, phenomena possess symmetry.} To take some one phenomenon as central is, before the analysis has even begun, to perform upon the phenomenal layer a \emb{spontaneous symmetry breaking}. Prior to analysis, gift, rite, address, gaze, absence, triangle---the phenomena co-exist symmetrically: none is \emph{a priori} more ``central'' than another. The gift comes to \emc{seem} central only because our attention happens to congeal, to break, there. That breaking is real, but it is an event on the \emb{observer's} side, not a structure on the side of the real. To fix the contingent landing-point of an observer's attention as a hierarchy in the theory is to mistake an epistemic contingency for an ontological necessity. The ``importance'' of the gift, if real, can only be that which its \emb{position} in the external relational field confers upon it (it happens to be richly connected), not some essence intrinsic to it. Load-bearing is conferred by relation, not carried by the node.

\emc{Second, no phenomenon contains infinite reason.} To take a single phenomenon as one's lens carries an over-strong assumption: that from it the entire theoretical field can be mirrored, as though it contained every dimension. But any phenomenon is finite, local. It illuminates only those dimensions strongly coupled to it and is \emb{blind} to the dimensions orthogonal to it. The gift illuminates the circulation of value and semantic coding, yet it cannot illuminate a silent rite in which nothing is exchanged (the rising and yielding of one's seat before an elder), nor a pure gaze attached to no object (the eyes that appraise one again and again across the banquet table). To force the gift to explain these phenomena it was never able to reach only distorts them into variants of the gift, losing the distinct phenomenological texture each one has. To let a local usurp the global is an epistemic overreach. This accords with a vigilance constant to the series: no finite thing---be it a subject, a relation, or a phenomenon---is entitled to the view from nowhere.

\emc{Third, relational being demands organization by coupling rather than by a hub.} This is the deepest of the three. The core ontological claim of this series is that load-bearing, meaning, and value are not intrinsic properties of any single node but \emb{emergent effects of the relational network}. Hub-and-spoke organization---one central phenomenon with the rest arrayed around it---violates precisely this claim, for it loads the weight of organization onto a single node and lets the network degenerate into a star centered on that node. The mode of organization faithful to the relational ontology is to let \emb{the coupling relations themselves} be the sole organizer: the phenomena hang on the relational network as equals, organized by \emb{the edges among them} rather than radiating from some specially chosen point. In such an organization the gift is merely a node of higher connectivity in the network---a connectivity that will be \emb{displayed} in the coupling matrix of \xref{sec:matrix}, not \emb{presupposed}. The difference between displaying and presupposing is precisely the difference between fidelity and infidelity to the relational ontology.

\emc{From this follows an original methodological claim of the paper: formal self-reference.} A paper that expounds ``how relational being unfolds toward the world'' ought itself to have a chapter-structure that is relational, decentered, organized by coupling rather than by a hub. In other words, the \emb{form} of this paper should self-referentially enact its \emb{content}. The earlier papers already carried this aesthetic awareness (\pinkref{Paper~III} wrote self-legislation in the form of self-legislation; the mother paper expounded relationality in a relational style); the present paper makes it thematic: the whole is organized not around a central phenomenon but along a vertical axis of ``how coupling occurs, layer by layer'' and a horizontal axis of ``how medium and tension cut across the layers,'' so that in the very experience of reading, the reader lives through one unfolding of relation, from the dyad toward the world.

One objection must be met in advance: does decentering amount to having no focus, issuing in diffuseness? It does not. Decentered organization has its definite spine---it is doubly oriented by the \emb{order in which coupling occurs} (from the first other beyond the dyad, to network and field) and by the \emb{theoretical strata} (the Real, the Symbolic, the phenomenal; see the stratal coordinates of \xref{sec:method}). The focus lies not in any phenomenon but in \emb{structure}: in how coupling ascends, in how the strata mutually ground one another. This is exactly the form a focus should take in a synoptic mother-text---it offers terrain and coordinates, not a close-up of any one vista.

% ============================================================
\section{Theoretical Framework: The Lenses of the Survey}\label{sec:frame}

A programmatic review earns its standing as infrastructure chiefly here, in the framework section, for it is here that a reader new to the domain is given the traditions to command and the literature to read. This section therefore does not stint. Each of the twelve frameworks the paper employs is given its own exposition and its own literature-review table, so that a researcher entering by any one lens finds, in one place, what the framework claims, who its principal authors are, what each contributes, and where the framework's force and its limits lie. The substantive \emph{application} of each framework still unfolds in the stratal chapter where it does its work---that is the principle of \pinkref{§2}, that the frameworks are nodes given sense by their couplings---but the \emph{exposition} belongs here, gathered, so that the later chapters may use the lenses without stopping to introduce them.

A word on the literature tables. They are reviews, not bibliographies: each entry states not merely a citation but the work's core contribution, its specific significance for this research programme, and its principal limitation or the controversy attaching to it. Every entry has been bibliographically verified; an accurate table is worth more than a comprehensive-looking but unreliable one, and the reader is owed the distinction. The frameworks are presented in the order of the strata they will serve---ontological ground first, then the symbolic frameworks (value, sign, law, economy), then the mechanism frameworks (cognition, neuroscience, decision), then the formal and the psychoanalytic---and the section closes (\pinkref{§3.12}) on the claim that binds them: that they form a network with no foundation, the third place the paper's form enacts its content.

\subsection{The ontology of relational being}
This is the framework carried from the first six papers and the ground on which the present one builds, though---by its own logic---a ground that is itself a node rather than a foundation. Its claim is that being-in-isolation, the self set apart and only subsequently entering relations, is an abstraction; the concrete truth is that subjects are constituted in and through relation, generated by the bonds that a prior metaphysics treated as external accidents befalling an already-complete substance. The mother paper argued this for the dyad; the present paper extends it outward, so that diplomacy becomes the name for relational being \emph{unfolding toward the world}---the dyad, having generated itself in second-order coupling, now generating and being generated by the field beyond it.

The lineage of the view is worth setting out, because the paper's originality lies less in the bare thesis (which has antecedents) than in its formalization and its extension to higher order. The relational thesis has roots in Hegel's account of self-consciousness as achieved only through recognition by another \citep{hegel1807}; in Buber's distinction of the I--Thou relation from the I--It \citep{buber1923}; in Marx's sixth thesis on Feuerbach, that the human essence is the ensemble of social relations \citep{marx1845theses}; and, in a different key, in the Buddhist doctrine of dependent origination (\zh{緣起}, \emph{pratītyasamutpāda}), that no entity possesses own-being but arises only in dependence on conditions. Contemporary relational sociology \citep{emirbayer1997manifesto,emirbayer1998agency,crossley2011towards,donati2011relational} and the relational turn in feminist ethics \citep{mackenzie2000relational} are its nearer kin. Against this background the series' contribution is to treat relation not as a fact about persons but as a generative ontology with formal structure---which is what permits the coupling-theoretic and geometric formalizations the later frameworks supply.

\goldrule
\footnotesize
\begin{longtable}{Y Y Y Y}
\caption{Relational being: literature}\label{tab:lit-relbeing}\\
\toprule
\rowcolor{pageblush}
Source & Core contribution & Significance for the programme & Limitation / controversy \\
\midrule
\endfirsthead
\rowcolor{pageblush}
Source & Core contribution & Significance for the programme & Limitation / controversy \\
\midrule
\endhead
Hegel, \emph{Phenomenology of Spirit} (1807), master--slave dialectic & Self-consciousness achieved through recognition & The deep source of constitution-in-relation & Idealist frame; difficult \\
Buber, \emph{I and Thou} (1923) & The I--Thou vs.\ I--It relation & Relation as primary, not derivative & Aphoristic, hard to operationalize \\
Marx, \emph{Theses on Feuerbach} (1845), VI & Essence as the ensemble of social relations & Relation prior to the isolated individual & Compressed; much-contested reading \\
Dependent origination (\emph{pratītyasamutpāda}); Nāgārjuna & No own-being; arising-in-dependence & Non-Western root of the relational thesis & Doctrinal variants; translation \\
Mackenzie \& Stoljar (eds.), \emph{Relational Autonomy} (2000) & Autonomy as relationally constituted & Bridges relation and the normative & Anthology; uneven \\
Donati, \emph{Relational Sociology} (2011) & Society as made of relations & Contemporary relational social theory & Heavy apparatus \\
\citet{emirbayer1997manifesto}, ``Manifesto for a Relational Sociology'' & Substances vs.\ processes; relations as primary & The programmatic statement of the relational turn & Manifesto, not a worked theory \\
Crossley, \emph{Towards Relational Sociology} (2011) & Social worlds as networks of iterated interaction & Relation as lived trajectory & Network bias; light on inner constitution \\
The mother paper of this series & Generative relational being, formalized & The immediate ground of this paper & Self-citation; programme in progress \\
\bottomrule
\end{longtable}

\subsection{Generative justice and the circulation of value}
The second framework is the interface with Ron Eglash's generative justice \citep{eglash2016intro}, and it supplies the normative spine of the paper's political economy (\xref{sec:polecon}). Its guiding question is whether value, once generated within a relational network, \emph{circulates} within it---is retained by its producers, flows through the network and returns to those whose labor or care created it---or is instead \emph{extracted}, siphoned in one direction to a point of accumulation outside the network. Generative justice reframes justice as a question of circulation rather than of static distribution: the issue is not only who holds how much at an instant, but whether value is kept in motion among those who make it. Eglash develops the notion across ecological, labor, and expressive value, drawing on cybernetic and indigenous-knowledge traditions \citep{eglash2014garvey,eglash2017engineering,eglash2024computational}, and the present paper proposes the intimate relational field as a further, and ancient, domain of the same dynamic: diplomacy is among the oldest forms in which value circulates rather than being extracted, and equally a site where extraction wears the mask of reciprocity.

The framework couples tightly to two others. It supplies the normative reading of the geometric-phase formalism (\pinkref{§3.5}): extraction is, formally, the non-conservation of value around a circuit of an asymmetric field. And it inherits from the older critique of alienation---Marx's account of the worker separated from the product of labor---which it generalizes from waged labor to value of every kind, including the relational and the affective. The foundational statement is \citet{eglash2016intro}; the framework is elaborated across open-source and artisanal value \citep{eglash2014garvey,eglash2017engineering} and, most recently, in computational reparations and decolonial circular value flow \citep{eglash2024computational}.

\goldrule
\footnotesize
\begin{longtable}{Y Y Y Y}
\caption{Generative justice and value circulation: literature}\label{tab:lit-genjust}\\
\toprule
\rowcolor{pageblush}
Source & Core contribution & Significance for the programme & Limitation / controversy \\
\midrule
\endfirsthead
\rowcolor{pageblush}
Source & Core contribution & Significance for the programme & Limitation / controversy \\
\midrule
\endhead
\citet{eglash2016intro}, ``An Introduction to Generative Justice'' & Value circulation vs.\ extraction; the canonical definition & The paper's normative spine & Normative baseline still developing \\
\citet{eglash2024computational}, ``Computational Reparations as Generative Justice'' & Unalienated circular value flow; decolonial transitions & The recent major statement & Scope of ``generative'' contested \\
Marx, \emph{Economic and Philosophic Manuscripts} (1844) & Alienation: worker severed from product & The ancestor of the extraction concept & Early Marx; reading debated \\
Mauss, \emph{The Gift} (1925) & The gift kept in circulation; \emph{hau} & Circulation as an ancient relational form & Archaic-case generalization \\
Ostrom, \emph{Governing the Commons} (1990) & Commons sustained without extraction & Institutional kin of circulation & Not framed as ``value'' per se \\
\bottomrule
\end{longtable}

\subsection{Semiotics, the semantics of the sign, and hermeneutics}
The third framework concerns how value and the rites are generated, coded, and read as signs, and it serves the chapter on the rites and value (\xref{sec:rite}). Its claim is that a gift, a form of address, a seating, an absence is a \emph{signifier} whose relation to its signified---affection, recognition, standing, refusal---is conventional, culturally specific, and learnable, so that the same object can bear opposite meanings across two grammars (the white flowers that honor in one code and mourn in another). Three strands are needed. \emph{Structural semiotics} \citep{saussure1916} gives the signifier/signified relation and the arbitrariness of the sign; \emph{Peircean semiotics} \citep{peirce_cp} adds the triad of icon, index, symbol and the interpretant, which matters because a gift signifies partly by resemblance and partly by causal connection, not by convention alone; and \emph{hermeneutics} \citep{ricoeur1976,gadamer1960} gives the theory of \emph{reading}, for the receiver of a sign must interpret it, and interpretation is fallible, situated, and---as \xref{sec:power} insists---unequally authorized. Misreading a sign is the diplomatic incident the phenomenal chapter (\xref{sec:network}) analyzes; the right to fix a sign's reading is the hermeneutical injustice the justice chapter (\xref{sec:justice}) anatomizes.

The framework's nearest application to gifts specifically runs through Baudrillard's sign-value \citep{baudrillard1972} and through the anthropology of exchange, but its philosophical weight is hermeneutic: it establishes that value in the relational field is not read off objects but \emph{conferred in interpretation}, which is exactly why the authority to interpret becomes a site of power.

\goldrule
\footnotesize
\begin{longtable}{Y Y Y Y}
\caption{Semiotics, semantics, and hermeneutics: literature}\label{tab:lit-semiotics}\\
\toprule
\rowcolor{pageblush}
Source & Core contribution & Significance for the programme & Limitation / controversy \\
\midrule
\endfirsthead
\rowcolor{pageblush}
Source & Core contribution & Significance for the programme & Limitation / controversy \\
\midrule
\endhead
Saussure, \emph{Course in General Linguistics} (1916) & Signifier/signified; arbitrariness & The sign-structure of the rites & Langue/parole dualism contested \\
Peirce, \emph{Collected Papers} (icon/index/symbol) & The interpretant; sign triad & Gifts signify by more than convention & Notoriously systematic, dense \\
Ricoeur, \emph{Interpretation Theory} (1976) & Reading as the recovery of meaning & Interpretation of the diplomatic sign & Broad; many works to triangulate \\
Gadamer, \emph{Truth and Method} (1960) & The situatedness of understanding & Why readings are positioned, fallible & Long; hermeneutic-circle worries \\
Baudrillard, \emph{For a Critique of the Political Economy of the Sign} (1972) & Sign-value distinct from use/exchange & The third value before relational value & Polemical; later turn contested \\
\bottomrule
\end{longtable}

\subsection{Optimization and dynamical systems}
The fourth framework is the first of the formal layer, and it serves the behavioral reading of the network chapter (\xref{sec:network}). It makes two proposals. The first is that diplomacy is a \emph{multi-objective constrained optimization}: the dyad must satisfy, at once, internal accord, the expectations of each family of origin, ambient social norms, and its own values---objectives that frequently conflict, so that there is no single act maximizing all, only trade-offs along a Pareto frontier. ``Decorum,'' on this reading, is not a code obeyed but the \emph{feasible set}: the region of acts violating no hard constraint, within which the diplomat must still choose. The second proposal is that the relational field is a \emph{dynamical system}: gifts and slights are perturbations, debt and reciprocity are feedback loops, and the field has stable basins (a relation that absorbs ordinary shocks) and failure modes (collapse into conflict, or rigidification into hollow form). The art of diplomacy is then keeping the system in a \emph{living} stable basin---deep enough to absorb perturbation, shallow enough that feeling still flows.

The framework is offered as a structuring vocabulary, not a predictive model; its risk, named honestly, is functionalist drift---mistaking a suggestive formalism for an established mechanism---and the paper guards against this by leaving worked models to the dedicated computational sub-paper and using the vocabulary here only to make ``decorum'' and ``tact'' precise rather than to predict behavior.

\goldrule
\footnotesize
\begin{longtable}{Y Y Y Y}
\caption{Optimization and dynamical systems: literature}\label{tab:lit-optim}\\
\toprule
\rowcolor{pageblush}
Source & Core contribution & Significance for the programme & Limitation / controversy \\
\midrule
\endfirsthead
\rowcolor{pageblush}
Source & Core contribution & Significance for the programme & Limitation / controversy \\
\midrule
\endhead
Pareto, multi-objective optimality (Pareto front) & Trade-offs without a single optimum & ``Decorum'' as the feasible/efficient set & Silent on how to choose among optima \\
Boyd \& Vandenberghe, \emph{Convex Optimization} (2004) & Constrained optimization, rigorously & The formal frame for the feasible set & Assumes convexity rarely present \\
Strogatz, \emph{Nonlinear Dynamics and Chaos} (1994) & Attractors, basins, stability & The relational field as dynamical system & Application here is qualitative \\
\citet{gottman2002math}, \emph{The Mathematics of Marriage} & Nonlinear dynamical modeling of dyads, empirically & A precedent for relational dynamics & Dyadic, not higher-order; contested \\
\bottomrule
\end{longtable}

\subsection{Geometric phase, holonomy, and Value Foam}
The fifth framework is the paper's most distinctive formal proposal, and it serves the formal interfaces of \xref{sec:power} and \xref{sec:polecon}. Its intuition is geometric. When a vector is carried around a closed loop on a curved surface, ``always pointing straight ahead'' at each step, it returns to its starting point rotated; the rotation, the \emph{holonomy} of the loop, measures the curvature enclosed, and it is irreducible to any single step. The proposal is that a season of diplomacy is such a loop: the dyad sets out from a relational state, passes through a sequence of exchanges and occasions, and returns to what looks locally like its starting point, yet altered---and the alteration, like the rotation, is a property of the circuit as a whole, not of any transaction along it. Formally, with the relational state carried along a closed path $\gamma$ by a connection $A$, the holonomy
\[
  \mathrm{Hol}(\gamma)=\mathcal{P}\exp\!\oint_{\gamma} A
\]
is non-trivial, $\mathrm{Hol}(\gamma)\neq\mathbb{1}$, exactly when the circuit leaves the state irreversibly deformed. On an \emph{asymmetric} field the connection $A$ is itself asymmetric, so the deformation accumulated by the weaker party over one circuit differs from the stronger's: the geometric statement of unequal cost (\xref{sec:power}), and, read through generative justice, the non-conservation of value around the loop (\xref{sec:polecon}).

The framework draws on the geometric (Berry) phase of physics \citep{berry1984,simon1983}, on gauge theory's connection and holonomy, and on the series' own ``Value Foam'' formulation, which adapts these to value and alienation. Its honest limitation is twofold: a high abstraction barrier, and the burden of showing that the formalism does explanatory work rather than ornamenting the prose---a burden discharged not here but in the dedicated formal paper. \emph{The notation above is provisional and is to be aligned with the Value Foam formulation developed in the author's prior work in this series.}

\goldrule
\footnotesize
\begin{longtable}{Y Y Y Y}
\caption{Geometric phase, holonomy, Value Foam: literature}\label{tab:lit-holonomy}\\
\toprule
\rowcolor{pageblush}
Source & Core contribution & Significance for the programme & Limitation / controversy \\
\midrule
\endfirsthead
\rowcolor{pageblush}
Source & Core contribution & Significance for the programme & Limitation / controversy \\
\midrule
\endhead
Berry, ``Quantal phase factors\ldots'' (1984) & The geometric (Berry) phase & The physical root of the holonomy idea & Transposition to value needs argument \\
Nakahara, \emph{Geometry, Topology and Physics} (2003) & Connections, holonomy, fiber bundles & The mathematical apparatus & Heavy prerequisites \\
Simon (1983), holonomy interpretation of Berry phase & Holonomy reading of the phase & Names the circuit-invariant precisely & Technical \\
The author's Value Foam formulation (this series) & Holonomy applied to value and alienation & The immediate formal source & Programme in progress \\
Marx, alienation (cf.\ \pinkref{§3.2}) & Value severed from producer & The phenomenon the holonomy formalizes & Bridging to geometry is novel \\
\bottomrule
\end{longtable}

\subsection{Theory of mind and the joint-attentional field}
The sixth framework is the cognitive one, serving the mechanism bridge and the triad (\xref{sec:bridge}, \xref{sec:triad}). To know how an external other appraises the relation requires \emph{mentalizing}---representing another's mental states---and higher-order coupling requires this mentalizing to be \emph{nested}: I represent how my mother appraises my partner, and how my partner imagines that appraisal, and so on up. The framework supplies, first, the developmental and comparative science of theory of mind \citep{premack1978,wimmer1983} (the false-belief paradigm; the question of its phylogenetic distribution); second, the account of \emph{joint attention} \citep{tomasello2008}---two minds attending to the same object and each aware of the other's attending---which \pinkref{Paper~V} used to characterize the dyad and which the present paper extends, since in diplomacy the dyad must either draw a third into joint attention or be drawn into a third-party-dominated attentional field. Whose attentional field governs a scene is, the paper argues, a question of power.

The framework's empirical core (developmental psychology, comparative cognition) is robust; its contested edge is the neural-localization claims that the next framework takes up, and the depth-of-nesting question---how many orders of mentalizing humans routinely sustain---which the bridge chapter treats as a power-laden variable rather than a fixed ceiling.

\goldrule
\footnotesize
\begin{longtable}{Y Y Y Y}
\caption{Theory of mind and joint attention: literature}\label{tab:lit-tom}\\
\toprule
\rowcolor{pageblush}
Source & Core contribution & Significance for the programme & Limitation / controversy \\
\midrule
\endfirsthead
\rowcolor{pageblush}
Source & Core contribution & Significance for the programme & Limitation / controversy \\
\midrule
\endhead
Premack \& Woodruff (1978), ``Does the chimpanzee have a theory of mind?'' & Coined the theory-of-mind question & The origin of the mentalizing frame & Animal case still debated \\
Wimmer \& Perner (1983), false belief & The experimental test of ToM & Operationalizes mental-state attribution & Task-format critiques \\
Tomasello, \emph{Origins of Human Communication} (2008) & Joint attention, shared intentionality & The basis \pinkref{Paper~V} extends & Strong human-uniqueness claims \\
Baron-Cohen, \emph{Mindblindness} (1995) & Mentalizing as a system & Nested mentalizing as machinery & Modularity contested \\
Dennett, the intentional stance; orders of intentionality & Recursive belief attribution & The formal sense of ``higher order'' & Stance-instrumentalism debated \\
\bottomrule
\end{longtable}

\subsection{Psychoanalysis}
The seventh framework supplies the account of the gaze, the triangle, transference, and the big Other, and it serves the triad chapter (\xref{sec:triad}). Its claim is that the subject is constituted in the field of the Other's desire and the Other's look: the dyad, entering the symbolic network of family and norm, is constituted as an object by the gaze of the \emph{big Other} \citep{lacan1964sem11}---the anonymous order of social expectation---and proposal and wedding are, in this register, symbolic registrations before that gaze. The framework also supplies the \emph{triangle}: the entry of a third party into a prior dyad activates a structure of mediation, jealousy, and exclusion (the in-law triangle), and the mechanism of \emph{transference}, by which an entering partner is made to carry affect that belongs to an older, unsymbolized history (the traumatic real of \xref{sec:real}). Sartre's analysis of the look \citep{sartre1943} and Lévinas's of the face \citep{levinas1961} give the two poles between which \xref{sec:triad} situates the diplomatic gaze.

The framework's power is its account of what eludes the rational-actor models of the bridge chapter---the carried, the unspoken, the transferential; its honest liability is the contested empirical testability of its core claims, which the paper handles by using psychoanalysis to \emph{describe structure} (the triangle, the gaze) rather than to make causal predictions.

\goldrule
\footnotesize
\begin{longtable}{Y Y Y Y}
\caption{Psychoanalysis: literature}\label{tab:lit-psych}\\
\toprule
\rowcolor{pageblush}
Source & Core contribution & Significance for the programme & Limitation / controversy \\
\midrule
\endfirsthead
\rowcolor{pageblush}
Source & Core contribution & Significance for the programme & Limitation / controversy \\
\midrule
\endhead
Lacan, \emph{Seminar XI} (1964), the gaze & The subject under the Other's look; \emph{objet petit a} & The mechanism of being-watched-as-whole & Textually difficult; school disputes \\
Lacan, the big Other & The symbolic order as Other & The anonymous gaze of norm & Contested formalization \\
Sartre, \emph{Being and Nothingness} (1943), ``the look'' & Objectification, shame, being-for-others & The objectifying pole of the gaze & Pessimism of for-others \\
Freud, on transference & Affect displaced onto a present figure & The entering partner as carrier & Clinical origin; generalization \\
Bowen, family systems; triangulation & The triangle as a stabilizing structure & The in-law triangle, formalized & Empirical basis contested \\
\bottomrule
\end{longtable}

\subsection{Phenomenology: intersubjectivity, co-presence, and the gaze}
The eighth framework is phenomenology, and it deserves independent standing rather than absorption into psychoanalysis, because one of the paper's three governing dimensions---the gaze of the Other, named in the subtitle---is at bottom a phenomenological matter, as is the whole stratum of co-presence, embodied being-with, and the constitution of a shared world that the diplomatic field presupposes. Phenomenology asks how the other is given to me at all, how a ``we'' is constituted in lived experience, and how being-seen transforms the seen---questions prior to, and presupposed by, both the psychoanalytic account of the gaze and the cognitive account of mentalizing.

Four strands are needed. \emph{The constitution of the other} (Husserl's fifth \emph{Cartesian Meditation}): the other is given not as an object inferred but as an \emph{alter ego} apprehended through the analogical appresentation of a body like my own---the founding problem of how intersubjectivity is possible at all, and the one the diplomatic field inherits when a stranger must be received as a subject. \emph{Embodied intersubjectivity} (Merleau-Ponty): the \emph{intercorporeality} by which bodies read one another beneath explicit thought, so that the diplomatic field is navigated first by the body---the stiffening, the averted eye, the held hand---before any proposition is formed. \emph{The social world and the we-relation} (Schutz): the phenomenological sociology of how a shared world is built from the face-to-face ``we-relation,'' the most direct interface to the paper's central object, the dyad as a ``we'' that must now extend its we-relation to a third. And \emph{the gaze} (Sartre, and Husserl's account of being-an-object-for-another): how being-seen converts the subject into an object in another's world, the structure §\ref{sec:triad} develops as being-watched-as-a-whole.

Crucially, phenomenology is not only a classical resource but a living and critical one, and the recent \emph{critical phenomenology} movement is what lets the framework carry the paper's attention to power. Where classical phenomenology described the structures of experience as if neutral, critical phenomenology---Ahmed's \emph{Queer Phenomenology}, with its analysis of \emph{orientation} (how bodies are directed, what is within reach, who is ``in line''); Guenther's account of how social structures deform the very capacity for experience; Ortega's analysis of multiplicitous selfhood---shows those structures to be shaped by power, race, gender, and history. This is the strand that couples phenomenology to §\ref{sec:power}: the gaze is not a neutral structure but one that falls along a gradient, and orientation is not freely chosen but inherited and enforced. The framework thus spans the classical (how the other and the we are constituted) and the critical (how that constitution is bent by power), which is exactly the doubled vision the paper needs for the gaze of the Other.

\goldrule
\footnotesize
\begin{longtable}{Y Y Y Y}
\caption{Phenomenology: literature}\label{tab:lit-phenom-frame}\\
\toprule
\rowcolor{pageblush}
Source & Core contribution & Significance for the programme & Limitation / controversy \\
\midrule
\endfirsthead
\rowcolor{pageblush}
Source & Core contribution & Significance for the programme & Limitation / controversy \\
\midrule
\endhead
\citet{husserl1931cm}, fifth \emph{Cartesian Meditation} & The other as \emph{alter ego}; analogical appresentation & How the other is given at all & Charge of residual solipsism \\
\citet{merleauponty1945} & Embodied perception; intercorporeality & The body as first organ of diplomacy & Dense; pre-thematic claims hard to test \\
\citet{schutz1932} & Phenomenological sociology; the we-relation & The dyad as a lived ``we'' extended to a third & Tension with Husserl's transcendental aim \\
Sartre, \emph{Being and Nothingness} (1943), ``the look'' & Being-seen as objectification & The objectifying pole of the gaze & Pessimism of being-for-others \\
Heidegger, \emph{Being and Time} (1927), \emph{Mitsein} & Being-with as a basic structure of existence & Co-presence as constitutive, not added & Obscurity; political controversy \\
\citet{ahmed2006queer} & Orientation; what is ``in line'' and within reach & Couples the gaze to power and inheritance & Essayistic; selective with sources \\
\citet{guenther2013} & Critical phenomenology; structures that deform experience & The gaze as power-laden, not neutral & Extends from an extreme case \\
\citet{ortega2016} & Multiplicitous selfhood; being-between worlds & The cross-cultural dyad's doubled self & Specific to Latina feminist context \\
\citet{zahavi2014self} & Contemporary synthesis: self, empathy, shame & The current state of intersubjectivity theory & Sides in internal debates \\
\bottomrule
\end{longtable}

\subsection{Jurisprudence}
The eighth framework, serving \xref{sec:law}, treats the external relational field as structured at once by \emph{informal} norms (the rites) and \emph{formal} law (marriage, inheritance, maintenance, the legal characterization of brideprice, the standing of in-laws, divorce, domestic violence). If the rites are the soft grammar, law is the hard; the two coexist, strain against one another, and convert---a brideprice given as a rite becomes, when an engagement breaks, a debt the courts will hear. The framework supplies \emph{legal pluralism} \citep{griffiths1986,merry1988,moore1973semiautonomous} (the coexistence and conflict of multiple normative orders: state law, lineage quasi-law, and, for cross-cultural dyads, two jurisdictions at once), including its application specifically to the domain of intimacy \citep{cohen2012intimacy} and the critical observation that legal default rules are not neutral but historically gendered, so that law hard-codes into enforceable form the asymmetries \xref{sec:power} traces in the soft grammar. It connects to the author's prior work on sequential legal judgment (the NeSy-POMDP framework), under which the legal characterization of an act is itself a partially observable sequential inference.

\goldrule
\footnotesize
\begin{longtable}{Y Y Y Y}
\caption{Jurisprudence: literature}\label{tab:lit-juris}\\
\toprule
\rowcolor{pageblush}
Source & Core contribution & Significance for the programme & Limitation / controversy \\
\midrule
\endfirsthead
\rowcolor{pageblush}
Source & Core contribution & Significance for the programme & Limitation / controversy \\
\midrule
\endhead
J.\ Griffiths, ``What Is Legal Pluralism?'' (1986) & Multiple legal orders coexisting & Rite/law/quasi-law as layered orders & ``Law'' arguably over-extended \\
S.\ E.\ Merry, ``Legal Pluralism'' (1988) & Survey of plural-order scholarship & Maps the field for the dyad & Descriptive breadth over depth \\
\citet{cohen2012intimacy}, legal pluralism in the domain of intimacy & Plural normative orders in intimate life & Law and rite in the same field & Focus on Western constitutional law \\
\citet{moore1973semiautonomous}, the semi-autonomous social field & How non-state fields generate binding norms & The rite as a quasi-legal order & Boundary of ``law'' contested \\
Pateman, \emph{The Sexual Contract} (1988) & The gendered substructure of contract & Law's gendered defaults & Strong thesis, contested \\
\bottomrule
\end{longtable}

\subsection{Political economy}
The ninth framework, serving \xref{sec:polecon}, asks what value diplomatic commerce \emph{creates} (relational, social, and symbolic capital; trust; recognition), how it \emph{circulates} (chains of reciprocity, the flow of the gift, the spread of reputation), and how it is \emph{extracted} (gendered emotional labor, the harvest of the superior party). It is the positive battlefield of generative justice (\pinkref{§3.2}), and it draws together Marx's value theory \citep{marx1867capital}, Mauss's account of the gift economy \citep{mauss1925gift} and its elaboration in economic anthropology \citep{levistrauss1949,sahlins1972}, Bourdieu's forms of capital \citep{bourdieu1986forms} (social and symbolic capital as convertible, non-monetary assets), and social-reproduction theory's analysis of unwaged care labor \citep{federici2004}. Its load-bearing claim for the paper is that relational value is genuinely \emph{produced}, not merely transferred---so that the question of just \emph{circulation} can even be posed---and its critical claim is that on an asymmetric field this produced value leaks to the powerful.

\goldrule
\footnotesize
\begin{longtable}{Y Y Y Y}
\caption{Political economy: literature}\label{tab:lit-polecon}\\
\toprule
\rowcolor{pageblush}
Source & Core contribution & Significance for the programme & Limitation / controversy \\
\midrule
\endfirsthead
\rowcolor{pageblush}
Source & Core contribution & Significance for the programme & Limitation / controversy \\
\midrule
\endhead
Marx, \emph{Capital} I (1867); value theory & Use/exchange value; surplus; alienation & The relational fourth value; extraction & Labor theory of value contested \\
Mauss, \emph{The Gift} (1925) & The three obligations; \emph{hau} & The gift economy's micro-justice & Archaic generalization \\
Bourdieu, ``The Forms of Capital'' (1986) & Social and symbolic capital & Relational capital as convertible & Economism charge \\
Polanyi, \emph{The Great Transformation} (1944) & Embeddedness; substantive economy & Diplomacy as embedded exchange & Historiographical disputes \\
Federici, \emph{Caliban and the Witch} (2004); social reproduction & Unwaged reproductive/care labor & The gendered structure of extraction & Scope of ``reproduction'' debated \\
\bottomrule
\end{longtable}

\subsection{Neuroscience}
The tenth framework, serving the bridge chapter (\xref{sec:bridge}), offers a cautious \emph{implementational} account---with no reductionist commitment---of how the relational state is inferred. The relational state is treated as a latent variable on which the subject performs approximate Bayesian inference (active inference, in Friston's free-energy formulation \citep{friston2010fep,friston2017process,parr2022active}), realized in candidate substrates: the medial prefrontal cortex in mentalizing and self/other value; the temporo-parietal junction in belief attribution; the orbitofrontal and ventromedial prefrontal cortex in value and social-reward coding; the dorsolateral prefrontal cortex in norm-compliance and fairness decision (the ultimatum-game literature). The framework's place in the paper is guarded twice over (\xref{sec:bridge}): against \emph{reduction} (implementation is not identity; the value of a recognition is not exhausted by its neural realization) and against \emph{mere metaphor} (the formalism must yield predictions, an explanation, and a definition, or it is decoration). It connects to the author's computational-neuroscientific work (EEG and grammar-based modeling at NCNP), which shares the latent-state/noisy-observation/sequential-update posture.

\goldrule
\footnotesize
\begin{longtable}{Y Y Y Y}
\caption{Neuroscience: literature}\label{tab:lit-neuro}\\
\toprule
\rowcolor{pageblush}
Source & Core contribution & Significance for the programme & Limitation / controversy \\
\midrule
\endfirsthead
\rowcolor{pageblush}
Source & Core contribution & Significance for the programme & Limitation / controversy \\
\midrule
\endhead
Friston, ``The free-energy principle'' (2010) & Active inference; the Bayesian brain & The unifying frame for state inference & Falsifiability much debated \\
Amodio \& Frith (2006), mPFC \& social cognition & Medial PFC in mentalizing/value & Substrate of higher-order mentalizing & Reverse-inference caution \\
Saxe \& Kanwisher (2003), TPJ \& belief attribution & TPJ in theory-of-mind tasks & Belief attribution's substrate & Localization replicability \\
Rangel, Camerer \& Montague (2008), value coding & OFC/vmPFC in valuation & Neural representation of value & Reductionism risk \\
Sanfey et al.\ (2003), ultimatum-game fMRI & dlPFC/insula in fairness decision & Substrate of norm-laden decision & Single-paradigm generalization \\
\bottomrule
\end{longtable}

\subsection{Decision theory}
The eleventh framework, also serving the bridge (\xref{sec:bridge}), closes the chain: once the relational state is inferred, the subject must \emph{decide}---what to give, whether to attend, how to answer an affront---under uncertainty. The natural formalism is sequential decision under partial observability, the POMDP \citep{kaelbling1998pomdp,sutton2018rl}, since one acts not on the state but on a belief about it (connecting to the author's NeSy-POMDP work on sequential judgment). Game-theoretic structure recurs: \emph{coordination} games (two families converging on one wedding form), \emph{signaling} games \citep{spence1973,zahavi1975} (the gift as a costly signal, credible because costly), and \emph{repeated} play \citep{axelrod1984} (return-gifts as reciprocal strategy). And decision proceeds under \emph{social preferences}---inequity aversion, a taste for reciprocity \citep{fehr1999}---rather than narrow self-interest, which is why a diplomatically rational agent pays to right an imbalance. The framework's idealizations (well-defined payoffs, common knowledge) are its limitation; its value is to make ``tact'' a precise notion---a policy minimizing expected surprise across all parties' inferences at once.

\goldrule
\footnotesize
\begin{longtable}{Y Y Y Y}
\caption{Decision theory: literature}\label{tab:lit-decision}\\
\toprule
\rowcolor{pageblush}
Source & Core contribution & Significance for the programme & Limitation / controversy \\
\midrule
\endfirsthead
\rowcolor{pageblush}
Source & Core contribution & Significance for the programme & Limitation / controversy \\
\midrule
\endhead
Sutton \& Barto, \emph{Reinforcement Learning} (1998/2018) & Sequential decision; value functions & The formal frame for diplomatic action & Assumes stationarity, reward \\
Kaelbling, Littman \& Cassandra (1998), POMDPs & Decision under partial observability & The relational state is partially observed & Computational intractability \\
Spence, ``Job Market Signaling'' (1973) & Costly signaling & The gift as a credible costly signal & Equilibrium-selection issues \\
Axelrod, \emph{The Evolution of Cooperation} (1984) & Reciprocity in repeated play & Return-gifts as reciprocal strategy & Idealized tournament setting \\
Fehr \& Schmidt (1999), inequity aversion & Social preferences in decision & Why agents pay to right imbalance & Parametric, contested fits \\
\bottomrule
\end{longtable}

\subsection{The coupling of the frameworks}
The twelve frameworks pull on one another, and the pulls are not optional cross-references but the very thing that gives each its sense in this paper. Semiotics (\pinkref{§3.3}) supplies the sign whose value political economy (\pinkref{§3.10}) then asks to circulate; optimization (\pinkref{§3.4}) formalizes the decision that decision theory (\pinkref{§3.12}) frames and neuroscience (\pinkref{§3.11}) implements; theory of mind (\pinkref{§3.6}) describes the mentalizing that the neural framework realizes and the triad enacts; psychoanalysis (\pinkref{§3.7}) and phenomenology (\pinkref{§3.8}) together articulate the gaze, the one as the structure of desire and the other as the structure of being-seen; jurisprudence (\pinkref{§3.9}) gives the hard grammar against which the rites (the semiotic and the relational frameworks together) are the soft; the geometric phase (\pinkref{§3.5}) supplies the formal skeleton on which generative justice (\pinkref{§3.2}) hangs its normative reading of extraction. Pull any one framework out and the others lose a connection, not merely a citation.

This is what it means to say there is \emb{no foundational framework beneath the rest}: one cannot order them as axioms and theorems, with relational being as the axiom and the rest derived, because relational being itself is given content only by what the other frameworks let it do in the world. It grounds them and is, in turn, specified by them---precisely the circularity a relational ontology embraces rather than flees. A foundation is a node that bears weight without being borne; there is no such node here. This is the third and final place the paper's form enacts its content: \pinkref{§2} argued it for phenomena (no central phenomenon), \xref{sec:triad} will show it for the triad (no master face among the three), and here it holds for the frameworks themselves (no master lens). A paper on relational being that rested on a foundational framework would refute itself in its architecture, however relational its sentences. \pinkref{Table~\ref{tab:frameworks}} gives the synoptic map of the twelve; it is to be read, like the coupling matrix of \xref{sec:matrix}, as a network of edges rather than a ranking of nodes.

\goldrule
\begin{landscape}
\footnotesize
\begin{longtable}{>{\bfseries\color{rosedeep}}p{2.8cm} Y Y Y}
\caption{The twelve frameworks: synoptic map}\label{tab:frameworks}\\
\toprule
\rowcolor{pageblush}
Framework & Core concept & Serves & Principal limitation \\
\midrule
\endfirsthead
\rowcolor{pageblush}
Framework & Core concept & Serves & Principal limitation \\
\midrule
\endhead
Relational being (\pinkref{§3.1}) & Constitution-in-relation & The whole; the ground & Circularity if read as foundation \\
Generative justice (\pinkref{§3.2}) & Circulation vs.\ extraction & \xref{sec:polecon} & Normative baseline developing \\
Semiotics \& hermeneutics (\pinkref{§3.3}) & Sign; reading; relational value & \xref{sec:rite} & Code-fixity overstated \\
Optimization \& dynamics (\pinkref{§3.4}) & Pareto set; attractor & \xref{sec:network} & Functionalist drift \\
Geometric phase / Value Foam (\pinkref{§3.5}) & Holonomy; phase difference & \xref{sec:power}, \xref{sec:polecon} & Abstraction barrier \\
Theory of mind (\pinkref{§3.6}) & Nested mentalizing; joint attention & \xref{sec:bridge}, \xref{sec:triad} & Localization debates \\
Psychoanalysis (\pinkref{§3.7}) & Gaze; triangle; big Other & \xref{sec:triad} & Testability contested \\
Phenomenology (\pinkref{§3.8}) & Intersubjectivity; co-presence; the gaze & \xref{sec:triad}, \xref{sec:network} & Classical neutrality vs.\ critical turn \\
Jurisprudence (\pinkref{§3.9}) & Formal vs.\ informal norms & \xref{sec:law} & Jurisdictional variance \\
Political economy (\pinkref{§3.10}) & Creation / circulation / extraction & \xref{sec:polecon} & Value metric contested \\
Neuroscience (\pinkref{§3.11}) & State inference; PFC & \xref{sec:bridge} & Reductionism risk \\
Decision theory (\pinkref{§3.12}) & POMDP; signaling; social prefs & \xref{sec:bridge} & Idealized assumptions \\
\bottomrule
\end{longtable}
\end{landscape}

% ============================================================
\section{Method and Apparatus: Review, Analysis, Mapping, Thick Description}\label{sec:method}

A programmatic review needs a method proper to it, different from the method of an ordinary argumentative paper. An argument marshals evidence toward a conclusion; a survey, by contrast, must \emph{lay out a field} so that others can work it---and its method is accordingly fourfold. \emb{Review} gathers the relevant traditions and states why each belongs (criteria of inclusion). \emb{Analysis} traces how the traditions and phenomena couple (the coupling analysis that is the paper's intellectual core). \emb{Mapping} renders the result navigable---the roadmap of sub-papers, the fields of the literature tables. And \emb{thick description} grounds the abstractions in recognizable scenes (the case studies). The four are ordered from the most general to the most particular, and together they discharge the paper's purpose, which is not to settle a question but to make a domain workable. The remainder of this section fixes the conventions of the two devices a reader will meet most often: the tables and the cases.

\subsection{The fields of the literature tables}
Each table carries four columns: the source (author, year, title); its core contribution; its \emb{significance for the programme} (why a later researcher must read it); and its limitation or the controversy around it. The tables are not a bibliography. A bibliography lists what has been read; these tables tell a future researcher what to read, in what order, and to what end---they are infrastructure, not record.

\subsection{The case-study apparatus}
The cases use \emb{composited, anonymized constructed instances}: each synthesizes recurrent situations rather than depicting a particular real person. This both protects privacy and, by composition, carries more explanatory force than a single anecdote, which is always hostage to its idiosyncrasies. The cases are concentrated in the two phenomenal chapters (\xref{sec:triad}, \xref{sec:network}); because the symbolic order has by then been laid out (\xref{sec:power}--\xref{sec:polecon}), the thick descriptions read not as isolated anecdotes but as enactments of structure. A case, in this paper, is never evidence \emph{for} a thesis; it is a site \emph{at which} the coupling of dimensions becomes visible. Nothing in the argument rests on the truth of any particular case, only on the recognizability of the structure it displays.

\subsection{The stratal coordinates}
The whole runs along a vertical geology: \emb{the Real} (the remainder that cannot be symbolized) $\rightarrow$ \emb{the Symbolic} (power / law / the rites / the machinery of value) $\rightarrow$ \emb{the mechanism bridge} (cognition and decision) $\rightarrow$ \emb{the phenomenal} (acts and thick description) $\rightarrow$ \emb{integration} (justice). The three names are taken from the topology of psychoanalysis, but they are used here for organization, not as a pledge of allegiance to any one school's full doctrine. The minimal commitment is only this: that acts in the phenomenal field presuppose a symbolic order that structures them, and that the symbolic order in turn fails to capture a real remainder. Power, being a structure of the Symbolic, must therefore appear \emb{before} the acts of the phenomenal field that it conditions---which is why the chapter on power (\xref{sec:power}) precedes the chapters on gift-giving and the rites, and not the reverse.

% ============================================================
\section{Surveying the Two Fields: The Dimensional Field and the Phenomenal Field}

This section gives the static map: the dimensional field, the phenomenal field, and the matrix of their coupling. The gift is, in that matrix, a node of \emb{high connectivity}; but its importance is an emergent of structure, not an intrinsic property---and here, in passing, the temptation to enthrone it as a hub is dissolved (in keeping with \S2). The map is static by design: it is the still photograph that precedes the moving account of how coupling occurs (\xref{sec:triad} onward).

\subsection{The dimensional field}
Three lenses, fixed by the paper's subtitle: value generation, justice, and the gaze of the Other. A word on why \emph{these} three, and not the twelve frameworks of \xref{sec:frame}. The frameworks are tools---disciplinary apparatuses one picks up. The dimensions are something else: the three \emph{questions} every diplomatic act answers whether or not anyone asks them. \emb{Value generation} is the question \emph{what is made and what is it worth}---and it subsumes semiotics, the semantics of the sign, hermeneutics, and optimization, all the ways value comes to be made, coded, read, and weighed. \emb{Justice} is the question \emph{is the making and sharing fair}---and it subsumes political economy, ethics, and the distributive--recognitive--rectificatory registers. \emb{The gaze of the Other} is the question \emph{under whose eyes does this occur}---and it subsumes the phenomenological and the psychoanalytic: the phenomenology of intersubjectivity and being-seen (how the other is given, how being-watched constitutes the watched), and the psychoanalysis of the gaze as the locus of desire and the big Other. The frameworks are how one studies; the dimensions are what is always already at stake. This is why three suffice where twelve were needed: the twelve are means, the three are the irreducible faces of the matter itself.

That the three are themselves coupled, not parallel, is the point that keeps this from being a tidy tripartite scheme. The gaze alters value (a gift's worth shifts under appraisal); value bears on justice (what is produced is what can be fairly or unfairly shared); justice inflects the gaze (the right to appraise is unequally held). One cannot treat one dimension and set the others aside, because each is partly constituted by its relations to the other two---the same relational structure the paper finds everywhere, here at the level of its own analytic categories. The burden of the later chapters is precisely to trace where the three cross, and the crossings are where the paper's real claims live.

\subsection{The phenomenal field}
Listed as equals, with no hierarchy implied by order: the gift; brideprice and dowry; the rites; forms of address; the first formal visit and its gift; hospitality; absence and presence; the triangle; seating and the order of toasts; cross-cultural coordination. Each is a node; none is the center. Some are object-laden (the gift), some are wordless (the yielding of a seat), some are purely spatial (seating), some purely temporal (the timing of a return gift). Their heterogeneity is the point: it is why no single one can serve as the lens for all.

\subsection{The coupling matrix}\label{sec:matrix}
What follows displays the figure itself rather than walking its edges one by one. \pinkref{Table~\ref{tab:matrix}} marks, for each phenomenon, which dimensions it strongly ($\bullet$) or weakly ($\circ$) activates. Read it not as a finished claim but as a conjecture-map: each cell is a hypothesis a later paper may confirm or revise. The gift's row is the most filled---it is the high-connectivity node. But the matrix \emb{shows} this; it does not \emb{presuppose} it. And it shows just as plainly that other rows light up dimensions the gift's row leaves dark---the row for address lights up power and culture where it barely touches value; the row for absence lights up justice and the gaze where it scarcely touches the semantics of the sign. This is the visual refutation of the hub. One further reading: the columns for \emb{power} and for \emb{cognition/decision} are nearly as full as the gift's row---a foreshadowing that power is a structure cutting across all phenomena (developed in \xref{sec:power}), and that the inference-and-decision bridge underlies all diplomatic acts (\xref{sec:bridge}).

\goldrule
\begin{landscape}
\footnotesize
\begin{longtable}{>{\bfseries\color{rosedeep}}p{3.4cm} c c c c c}
\caption{Phenomenon $\times$ dimension coupling matrix ($\bullet$ strong, $\circ$ weak)}\label{tab:matrix}\\
\toprule
\rowcolor{pageblush}
Phenomenon \textbackslash\ Dimension & Value & Justice & Gaze & Power & Cogn./Dec. \\
\midrule
\endfirsthead
\toprule
Phenomenon \textbackslash\ Dimension & Value & Justice & Gaze & Power & Cogn./Dec. \\
\midrule
\endhead
Gift / brideprice / dowry & $\bullet$ & $\bullet$ & $\bullet$ & $\bullet$ & $\bullet$ \\
The rites (\zh{礼}) & $\bullet$ & $\circ$ & $\circ$ & $\bullet$ & $\circ$ \\
Forms of address & $\circ$ & $\circ$ & $\circ$ & $\bullet$ & $\circ$ \\
Hospitality & $\bullet$ & $\circ$ & $\bullet$ & $\circ$ & $\circ$ \\
Absence \& presence & $\circ$ & $\bullet$ & $\bullet$ & $\bullet$ & $\circ$ \\
The triangle & $\circ$ & $\bullet$ & $\bullet$ & $\bullet$ & $\bullet$ \\
Seating \& toasts & $\circ$ & $\circ$ & $\bullet$ & $\bullet$ & $\circ$ \\
Cross-cultural coordination & $\bullet$ & $\bullet$ & $\circ$ & $\bullet$ & $\bullet$ \\
\bottomrule
\end{longtable}
\end{landscape}
\todo{the cell assignments are presently indicative; each warrants argument cell by cell in a later paper}

% ============================================================
\section{The Real: The Remainder That Cannot Be Symbolized}\label{sec:real}

Before entering the Symbolic---power, law, the rites, the machinery of value---one stratum must be marked that lies beneath it and that the Symbolic never wholly captures. This chapter is deliberately \emb{short and sharp}, and its brevity is principled, not a shortfall: the Real is, by its nature, what resists articulation, so a long chapter \emph{about} it would betray it, converting into discourse precisely what is defined by escaping discourse. The work here is not to develop the Real but to \emph{post} it---to mark its presence and its four faces---for a single structural reason. The chapters to come will mount a critique of symbolic power, and a critique that forgot the Real would slide into the totalizing reading \pinkref{§7} warns against: it would treat the Symbolic as the whole, power as omnipotent, the subject as fully constituted by the order that ranks her. Posting the Real first inoculates against this. It keeps open, beneath every symbolic structure however oppressive, a remainder the structure cannot reach---and that remainder, as the last face will show, is the ground of every possibility of change. The Real is marked here, then, not for its own sake but as the floor under the critique: the guarantee that domination, however total it appears, is never quite total.

\emc{Givenness.} One does not choose one's kin. A partner's parents, siblings, lineage are \emb{given}, not contracted; one enters a relational field whose terms were fixed before one arrived and cannot be renegotiated as a contract's terms can. This givenness is the first face of the Real in diplomacy: a brute facticity the rites may clothe but cannot dissolve. One may learn the correct address, master the order of toasts---and still the fact of \emph{whose} parents these are remains unchosen, unexchangeable, prior to all rite.

\emc{Inalienability.} Marcel Mauss saw in the gift a remainder that resists pure exchange: the \emph{hau}, the spirit of the thing given, which keeps the gift bound to its giver even after it has passed to another. The inalienable is what cannot be fully converted into exchange-value without violence to its nature. An heirloom, a parent's blessing, the body of the beloved: to price these is not to value them but to commit a category error felt as a wound. The Real is here the limit of commodification---the point at which the symbolic operation of pricing meets something it cannot, without damage, encode. This will return as the inner kernel of the recognition-versus-commodification problem in \xref{sec:justice}.

\emc{The traumatic real of the family.} Every family of origin carries a real that has never been symbolized: the unspoken loss, the wound around which the family's silences are organized, the affect that has no place in the family's official narrative. When a partner enters, they enter not only a symbolic order of positions but a field charged with this unsymbolized real, and they may, without knowing why, be made to carry it (the transferential mechanism of \xref{sec:triad}). Diplomacy here is not negotiation among clear positions but movement through a field whose deepest charges are precisely those that cannot be stated.

\emc{The remainder.} Not everything in the relational field can be ritualized. There is a debt that cannot be discharged by any return gift (the debt to a parent for one's existence), a gratitude with no adequate expression, a love that exceeds every sign offered for it. To mark this remainder is to keep open, against the critique to come, a space of \emb{the Real's resistance}: an oppressive symbolic structure, however total it appears, can never wholly subsume the real. This is not a consolation but a structural fact, and it is the ground of every possibility of change.

% ============================================================
\section{The Symbolic (I): The Topology of Power}\label{sec:power}

Power is not another phenomenon set beside the gift and the rites. It is a \emb{structure of the Symbolic}---the prior condition under which every ritual act of the phenomenal field comes to bear meaning, to be measured, to be judged. This is why the chapter on power must appear before the acts of gift-giving. To place it afterward, as an addendum, would be to let the phenomenal acts present themselves first and then patch in a structural layer that in truth underlies them. The order of exposition here tracks the order of grounding.

The central thesis of this chapter revises the diplomatic metaphor itself. The earlier sections, by default, let ``diplomacy'' connote negotiation among sovereign units, with its implied formal parity. But before parents and lineage, the dyad is \emb{often not a sovereign equal}. It enters a field of authority that exists prior to it, that is given rather than chosen, and that holds the power of recognition and the disposal of resources. The diplomacy here is closer to \emb{asymmetric diplomacy}---tribute, dependency, the patron-client bond of protection and fealty---than to a treaty between equals.

\subsection{Why diplomacy is so often asymmetric}
The diplomatic metaphor, taken naively, carries a presumption of parity: states treat as formal equals, however unequal in fact, and the very form of treaty presupposes two sovereigns who could, in principle, have declined. The first task of this chapter is to show that this presumption fails for the dyad's most consequential diplomacy---its dealings with parents and lineage---and to show that the failure is \emb{structural}, not incidental.

Begin with the descriptive resource and then press it. Fei Xiaotong's \emph{differential mode of association} (\zh{差序格局}, \emph{chaxu geju}) \citep{fei1947soil} figures the Chinese relational order as concentric ripples radiating from the self, obligation and intimacy declining with distance, every tie carrying an implicit rank. Read not as ethnography but as \emb{structure}, this is a topology of power: the positions in the order are not equal, and the position into which an in-marrying party is inserted is characteristically peripheral and low. Against the liberal image of marriage as a contract between two autonomous equals the contrast is sharp---where that image supposes parity, the differential order supposes rank, and the entering dyad is ranked on arrival.

Three objections must be met, for on their resolution the whole chapter turns.

\emc{First: is the asymmetry structural, or merely the contingent product of bad families?} One might grant that some in-laws are domineering while insisting that this is a defect of persons, not a feature of the field---that a generous family would receive the dyad as an equal. The reply is that the asymmetry survives the goodwill of every party. A family may wish to welcome the entering partner as an equal and still seat them according to the order, still hold the historical narrative the partner cannot access, still occupy the position from which appraisal flows rather than the position appraised. Goodwill can soften the exercise of the gradient; it cannot flatten the gradient itself, because the gradient is constituted by position---by seniority, by priority in time, by the direction in which recognition is sought---and position is not in anyone's gift to renounce. A kind sovereign is still a sovereign. That a structure can be inhabited gently is no evidence that it is not a structure.

\emc{Second: does }chaxu geju\emc{ over-generalize from one culture?} The differential mode is a Chinese model; to rest a general claim on it would be to commit, at the level of theory, the very cultural-capital error the chapter will name. The claim is therefore narrower and stronger than ``all relational fields have this shape.'' It is that \emb{wherever} a relational field assigns positions by seniority, priority, or the direction of recognition---and the anthropological record suggests this is near-universal, though its content varies enormously---the field has a power topology, and the in-marrying party enters it low. The differential mode is one richly described instance, valuable because it makes the topology explicit; it is offered as an illuminating case, not as the universal form. The Western contractual image is then not the culture-free baseline against which others deviate but \emph{another} local form---one that conceals its own gradients (of wealth, of family standing) behind the language of equal contract. No relational field is flat; cultures differ in whether they say so.

\emc{Third: does calling it ``tribute'' or ``fealty'' overstate the metaphor?} These words carry a charge of domination that may seem to caricature loving families. The point of retaining them is diagnostic, not pejorative. \emb{Asymmetric diplomacy} names a real and recognizable form---the lower-standing party petitions, defers, offers tokens of good faith, and seeks an acceptance that the higher party is free to withhold---and the affective warmth of a particular family does not change the form, any more than a benevolent suzerain changes the structure of tribute. To keep the hard word is to keep visible what gentler language would hide: that the entering dyad is, in this commerce, asking rather than treating, and that the power to grant or refuse lies on the other side.

Givenness completes the case for structurality. Because one cannot choose one's in-laws as one chooses a contracting partner, one cannot exit the asymmetry by the ordinary contractual means---declining the terms, walking from the table. The exits a contract affords are foreclosed; the relation is entered as a fact, not concluded as a bargain. The asymmetry is therefore structural in the strong sense: it is fixed by position, survives goodwill, and admits no contractual exit.

\subsection{The axes of power}
Power in the external relational field is not a single gradient but a layering of several, and the layering matters as much as the axes themselves. Take them first in turn.

\emb{The generational axis.} Seniority, and the duty of filial deference (\zh{孝}, \emph{xiào}), assign the junior a position of lower standing from which deviation reads as transgression. This axis is peculiar in that it is, in principle, temporary---the junior ages into seniority---which can make it feel less like domination than like a queue. But the temporariness is cold comfort to the party who is junior \emph{now}, and the in-marrying partner may remain junior, relative to the elders of a family not their own, for as long as those elders live.

\emb{The gendered axis.} Under patrilocal and patrilineal arrangements the in-marrying party---most often, historically, the wife---occupies a structurally weaker position, entering the husband's family as the one who must adapt, relocate, and prove. This axis, unlike the generational, does not dissolve with time; the wife who becomes a mother-in-law does not thereby cease to have entered low, and the structure reproduces itself in the next entering wife.

\emb{The economic axis.} Brideprice, housing, and the flow of money create dependencies that translate directly into standing: the party who receives owes, the party who owes defers. Money here does not merely buy goods; it buys position, and the debt it creates is paid in deference (the mechanism \xref{sec:justice} will read as the weaponization of reciprocity).

\emb{The axis of cultural capital.} Whose rites count as the ``default,'' whose mistakes are ``errors''---in cross-cultural dyads one side's culture is frequently naturalized as correct and the other's marked as deviation. This axis is the most invisible, because it operates through what \emph{needs no explanation}: the dominant grammar is simply ``how things are done,'' and only the other grammar requires justification.

\emb{The axis of information and the emotional ledger.} The family of origin holds the historical narrative and the long record of who owes whom what affect---an archive the in-marrying party can neither access nor contest. To be without this archive is to be perpetually one move behind: unable to know why a remark lands hard, which old wound a casual sentence reopens, what debt a gift silently answers.

\emc{The decisive point is the superposition.} These axes are not independent variables that might be assessed one at a time; they \emb{superpose}, and the same person may be lowered by several at once, with effects that compound rather than add. The young, in-marrying wife from the non-default culture, dependent on a brideprice her family received and ignorant of the husband-family's emotional ledger, is not lowered four times over in some additive sense; she is placed at the intersection where four gradients meet, and the intersection has a steepness no single axis predicts. This is why the analysis cannot proceed axis by axis and then sum: the structure is the superposition, and the superposition is what a later, dedicated paper on the power topology must formalize. One should resist, equally, the opposite error---treating the axes as a single undifferentiated ``oppression.'' They are distinct, they can pull against one another (an elder woman outranks a junior man on the generational axis while the gendered axis runs the other way), and a diplomacy that would navigate the field must read each, and read their interference.

\subsection{How power cuts across the dimensions}
The reason power earns a chapter in the Symbolic, rather than a place among the phenomena, is that it does not occupy one dimension but \emb{cuts across all of them}. A phenomenon sits at one or a few nodes of the coupling matrix; power runs through the whole column. Four cross-cuts are decisive, and each is the seed of a later chapter.

It is, first, \emb{the right to define value}. Who decides whether a gift is ``thick'' enough, whether an act is ``proper,'' whether a gesture counts as respect or as insolence? The standard is not neutral; it is held by the superior party, so that the very scale on which the dyad's offerings are weighed is calibrated by the side it must satisfy. This is the cross-cut into value generation (\xref{sec:rite}): value is not read off an object but conferred by whoever holds the measure.

It is, second, \emb{the right to interpret}---the diplomatic variant of the epistemic injustice of \pinkref{Paper~VI}. Who holds the authority to say what a gift, a silence, a tone of voice ``meant''? In the asymmetric field the superior party's reading stands as fact and the subordinate's as mere impression, to be corrected; the entering partner's account of a slight is received as oversensitivity, the family's account of the partner's conduct as simple description. To have one's interpretations systematically discounted is a wrong distinct from being given too little---it is hermeneutical, not distributive---and it is the cross-cut into justice (\xref{sec:justice}), where it compounds every other injustice by disabling the very capacity to contest them.

It is, third, the \emb{one-directionality of the gaze}. The elder appraises the junior, the husband's family appraises the new wife, and the converse does not hold; to gaze back, to appraise the appraisers aloud, is itself a breach. The gaze in the asymmetric field flows downhill, and the lower party is its object, rarely its subject. This is the cross-cut into the phenomenology of being watched (\xref{sec:triad}): whether the gaze objectifies or summons depends on the gradient down which it falls.

It is, fourth, ``decorum'' as the \emb{reproduction of power}. The standard of what is proper is set by those above; conformity to the rites is therefore, among other things, the continual re-enactment of the existing hierarchy. Every correct bow rehearses who bows to whom. This is the cross-cut into the rites as grammar (\xref{sec:rite}): the grammar is not neutral syntax but syntax whose rules encode rank.

\emc{One must guard, here, against a totalizing reading.} To say power cuts across all dimensions is not to say that everything is power, that recognition is only domination and the gift only a lien. That deflation would be as false as the naive parity it replaces, and it would make the chapter's own balanced standing incoherent. The claim is the more exact one: that power is a \emph{structuring} layer present in every dimension, not the \emph{whole} of any dimension. A gift is really a gift and also a move in a field of standing; an interpretation is really an understanding and also an exercise of the right to read. The rites are grammar---and the grammar's rules encode rank. Holding both halves of each ``and'' is the discipline this chapter demands; releasing either half falsifies the phenomenon.

\subsection{The formal interface: phase accumulation in an asymmetric field}
The cross-cuts just traced are qualitative; this section sketches how they might be made formal, and why the geometric-phase language of \xref{sec:frame} is the apt one. The aim is not to formalize here---that is a later paper's work---but to show that the chapter's central thesis, unequal cost, has a precise geometric statement, so that the deferral is a deferral of execution and not of conception.

Begin with the intuition. A season of diplomacy---an engagement, a wedding, the first year of visits---is a \emb{circuit}: the dyad sets out from some relational state, passes through a sequence of exchanges and occasions, and returns to what looks, locally, like where it began. The couple is married now as they were courting before; the visits resume their rhythm. Yet something has been accumulated that no single exchange records, the way a vector carried around a closed loop on a curved surface returns rotated though it was, at each step, carried ``straight.'' That accumulated rotation---the failure of the state to return truly to itself after a closed circuit---is what holonomy measures, and it is the natural formal home for the irreversible: the way a relation is permanently altered by a passage through the field, beyond any of the discrete transactions along the way.

The thesis of this chapter adds one thing to that picture: in an asymmetric field, \emb{the connection itself is asymmetric}. The ``connection,'' in this language, is what specifies how the relational state is carried from one occasion to the next; on a tilted field it carries the two parties differently. The weaker party, traversing the \emph{same nominal circuit}---the same wedding, the same visits, the same gifts given and returned---accumulates a different, and characteristically greater, deformation than the stronger. The objects exchanged may balance to the last token; the transformations undergone do not. This is the geometric statement of unequal cost, and it captures something the ledger of exchanges cannot: that two people can complete identical itineraries and be changed by them unequally, the lower party bent further by the passage than the higher. What \xref{sec:polecon} will describe as value leaking to the superior party, and \xref{sec:justice} as the weaponization of reciprocity, are, in this language, the same fact seen as non-conservation around an asymmetric loop. The full development---the specification of the connection, the conditions under which the holonomy is non-trivial, the measurement of the asymmetry---is left to the dedicated paper. \todo{align with Value Foam notation once confirmed}

\emc{A balanced standing.} It remains to forestall the misreading this chapter most invites: that to expose the reproduction of power through the rites is to indict the rites as such, to recommend their abolition. It is not. The rites are, as \xref{sec:rite} will argue, the formal cause by which relation takes shape---without them relation has no determinate contour---and they carry an ontological value and an irreplaceable face of care; the same bow that rehearses hierarchy also enacts a tenderness and a recognition that have no other vehicle. The position taken here is the harder, double one: \emb{to criticize the oppressive operation of the rites without denying the rites themselves}. The rites give relation its form, \emph{and}, in the hands of the powerful, that form reproduces domination---both at once, of the same rites, in the same act. The temptation is to collapse the double vision into one of its poles: into an uncritical reverence that hears only the tenderness, or a wholesale rejection that hears only the domination. Each pole is easier to hold and each falsifies the phenomenon, which is precisely that the rites are not two things, a good one and a bad one, but one thing that does both. The generative question---posed in \xref{sec:thirdrites} and left open---is whether a form can be found that keeps the care while loosening the domination, or whether the two are so bound in the existing rites that the only honest path is to make new ones.

\begin{casebox}{The structural vulnerability of the new wife}
A woman marries into a patrilocal household and, in the first months, finds her position fixed less by anything she does than by where the structure has placed her. The forms of address she must use ascend; those used toward her do not. At the table she is seated below, and serves before she is served. Her readings of slights are received as the over-sensitivity of an outsider; the family's readings of her conduct are received as fact. No single act here is cruel; the cruelty, such as it is, is structural---the superposition of the generational and gendered axes assigning her a low position from which her own interpretations carry little weight. The case activates, at once: the axes of power (gender and generation), the right of interpretation (\xref{sec:power}, and \pinkref{Paper~VI}), and the one-directionality of the gaze.
\end{casebox}

\begin{casebox}{Dependency manufactured by brideprice}
A large brideprice is paid; the marriage is thereby inflected, from its outset, by debt. What was given as a mark of seriousness functions, over time, as a lien: the receiving family's sense of obligation, the giving family's sense of entitlement, the couple's awareness that a price was named. Economic transfer converts, by degrees, into standing and dependency, and the party associated with the debt defers. The case activates: political economy, the justice of debt (\xref{sec:justice}), and the limit of commodification (\xref{sec:real})---the unease of brideprice lies precisely where a transfer meant to honor brushes against the pricing of a person.
\end{casebox}

\begin{casebox}{The marginalized rites of one side}
In a cross-cultural dyad, one partner's repertoire of rites is treated, by the wider field, as the default and correct one; the other's is treated as quaint at best, as error at worst, and the second partner is expected to assimilate. The asymmetry is rarely stated; it operates through the silent assignment of which practices need explanation and which do not. The case activates: the axis of cultural capital, and the right to define---here, the prior right to define whose culture is the unmarked ground. It anticipates the generative problem of the ``third set of rites'' (\xref{sec:rite}).
\end{casebox}

\goldrule
\footnotesize
\begin{longtable}{Y Y Y Y}
\caption{The topology of power: literature}\label{tab:power}\\
\toprule
\rowcolor{pageblush}
Source & Core contribution & Significance for the programme & Limitation / controversy \\
\midrule
\endfirsthead
\toprule
\rowcolor{pageblush}
Source & Core contribution & Significance for the programme & Limitation / controversy \\
\midrule
\endhead
Foucault, \emph{Discipline and Punish} (1975) & The microphysics of power & The rites as a micro-technique of discipline & Power as ubiquitous risks losing the agent \\
Bourdieu, \emph{Distinction} (1979); symbolic power & Cultural capital; symbolic violence & The class character of ``decorum'' & Determinism charge \\
Fei Xiaotong, \emph{From the Soil} (1947) & The differential mode of association & The native model of the power topology & Descriptive, pre-critical \\
Feminist care ethics \& family politics (e.g.\ Kittay, \emph{Love's Labor}) & The gendering of care & The extraction of emotional labor & Risk of essentializing care \\
Scott, \emph{Weapons of the Weak} (1985) & Everyday resistance of the subordinate & The agency of the weaker party & Romanticizing resistance \\
\bottomrule
\end{longtable}

% ============================================================
\section{The Symbolic (II): The Rites as Grammar, Value as Sign}\label{sec:rite}

If the previous chapter established that the field is tilted, this one asks what the field is \emph{made of}. The answer is that it is made of a grammar---the rites (\zh{礼}, \emph{lǐ})---within which alone the acts of diplomacy become meaningful, and within which value itself is generated as a kind of sign. The rites are the \emb{grammar} of the Symbolic: they fix positions and relations; they are the formal cause by which relation takes shape (carried over from the mother paper). If the rites are the soft grammar, law (\xref{sec:law}) is the hard one. And because the authority to interpret this grammar has already, in \xref{sec:power}, been distributed unequally along the gradient of power, this chapter must be read against the preceding one: the grammar is never innocent of the rank its rules encode. The order of treatment matters---grammar before utterance, as \xref{sec:method} insisted---because the phenomenal acts of the later chapters are utterances, and an utterance is unintelligible except in the grammar that makes it one.

\subsection{The rites as formal cause and symbolic grammar}
In the Confucian account, \emph{lǐ} is not an external constraint laid upon a relation that would exist without it; it is the form through which the relation comes to have a determinate shape at all. Without \emph{lǐ} the relation has no defined contour---no settled way for a son to stand before a father, for a guest to enter a host's house, for an in-marrying party to greet the lineage. The relation does not pre-exist its form and then receive it; it becomes the relation it is \emph{in} taking form, as the mother paper argued of relational being generally. This is the precise sense in which \emph{lǐ} is a formal cause and not a mere code of manners.

The claim that \emph{lǐ} is a \emph{grammar} must be defended against the charge that the word is borrowed loosely, a humanist's gesture at linguistics. The defense is that the analogy holds at the level that matters: \emph{lǐ} has the three properties that make a grammar a grammar rather than a list. It has a \emb{combinatorics}---occasions, persons, and acts compose by rules, so that the same gift is correct at one threshold and an affront at another, as the same word is grammatical in one position and not in another. It has a \emb{generativity}---competent participants produce and understand novel but well-formed acts they have never seen performed, the mark of a rule-governed system rather than a memorized repertoire; one knows, never having faced exactly this configuration of guests, what would be fitting. And it has \emb{grammaticality judgments}---participants agree, often without being able to state the rule, that a given act ``does not go,'' the felt wrongness of a solecism, which is the diagnostic of an internalized grammar. A system with combinatorics, generativity, and grammaticality is a grammar in more than name. This is also why \emph{lǐ} resists translation: ``rites'' captures the ceremonial, ``propriety'' the virtue-laden, ``etiquette'' the mannerly, but each loses the structural fact that \emph{lǐ} is the system within which relation becomes legible as the relation it is. The term is therefore kept in transliteration throughout. As grammar, \emph{lǐ} assigns positions; and to assign positions is already to do the work the previous chapter called power, which is why grammar and rank are, in the rites, inseparable.

\subsection{The semiotic generation of value and the semantics of the sign}
A gift is a sign, and its value is not its price. A hand-made thing may outweigh a costly one because the relation it signifies is denser; a returned heirloom may be priceless in a sense no sum captures. To see why, it helps to set out the series of value and then add to it. Classical political economy distinguishes \emph{use-value} (what a thing does) from \emph{exchange-value} (what it fetches); the semiotic tradition adds \emph{sign-value} (what a thing signifies about its owner, the Veblen-Baudrillard register in which a watch tells the time and also tells the world one's standing). To these the relational field adds a fourth, which this paper names \emb{relational value}: the value a thing has in that it \emph{generates or confirms a relation}.

That the fourth is genuinely distinct, not a relabeling of sign-value, can be shown by a case the others cannot accommodate. Consider a small, worthless object---a pressed flower, a borrowed book returned with a note---that is, between two people, of immense value, and of \emph{no} value to anyone else, and whose value would be \emph{destroyed} by sale. Sign-value cannot explain it: the object signifies nothing to the world, confers no standing, is invisible as a marker. Exchange-value is near zero and use-value negligible. Yet it is not valueless; it is, to the two, weighty. The weight is relational: the object's value consists in its having generated, and its continuing to confirm, a particular bond---and this value is, by its nature, \emb{non-transferable and non-priceable}, since to sell it is precisely to dissolve the relation that constituted its worth. Relational value is thus not high sign-value but value of another kind, indexed to a relation rather than to a market or an audience, and it is the value diplomacy chiefly trades in. (Its non-priceability is the same fact the Real named as inalienability in \xref{sec:real}, and the wound of its commodification is the subject of \xref{sec:justice}.)

The semantics of this sign-system is, like any grammar's, culture-specific and prone to drift. The same object signifies opposite things across cultures---a clock, an umbrella, white flowers carry, in Chinese contexts, associations of death, of separation, of mourning, that they do not carry elsewhere---so that a gift chosen for its evident loveliness can transmit its opposite. And the same transfer slides, with context, between gift and bribe: the boundary is not in the object or the sum but in the relation and the expectation of return, which is why the identical envelope is gratitude in one frame and corruption in another. To give well is to command this semantics; to misread it is the hermeneutic incident analyzed, at the level of the phenomenal act, in \xref{sec:network}. The point to carry forward is that value here is not read off the thing but \emph{conferred} in a sign-system whose code one party may command and another stumble in---which is exactly where the previous chapter's ``right to define value'' bites.

\subsection{Grammatical difference across cultures}
The grammar is not one, and the differences are differences of structure, not merely of vocabulary---which is why a dyad standing between two of them faces a problem deeper than translation. The Chinese differential rites assign elaborate, graded obligations by rank and distance, the concentric order of \xref{sec:power} made into a grammar of address, gift, and precedence. The Japanese economy of \emph{giri}--\emph{ninjō} (\zh{義理}--\zh{人情}) sets the weight of contracted obligation against the movement of spontaneous feeling, and governs the seasonal cycles of return-gift (\emph{ochūgen}, \emph{oseibo}) with a precision that can strike an outsider as an accountancy of the heart---though to read it as cold is itself a grammatical error, mistaking the careful discharge of \emph{giri} for the absence of \emph{ninjō} when the two are meant to be held together. Western \emph{etiquette} and \emph{hospitality} run flatter, more contractual, more nearly between equals, encoding rank less and reciprocity-among-peers more.

These are not three dialects of one grammar but three grammars, and the temptation to rank them---to treat one as more ``advanced,'' more egalitarian, less burdensome---should be resisted as itself an exercise of the cultural-capital axis (\xref{sec:power}). The flatter grammar is not the culture-free truth the others approximate; it encodes its own gradients (of class, of wealth, of family name) beneath the language of equal manners, and its very flatness can leave the newcomer without the explicit cues a graded grammar provides. Beneath the Western register, moreover, lies an ethics of the other that the philosophical tradition has made explicit and that the other grammars hold implicitly: Lévinas's face that commands before it is understood \citep{levinas1961}, the obligation to the guest that precedes any contract; and Derrida's aporia of hospitality \citep{derrida2000hospitality}---that a hospitality truly unconditional must welcome even the one who might destroy the home, so that the very rules which make hospitality possible (the host remains host, sets the terms) are what keep it from being unconditional. This aporia is not an exotic puzzle; it is the everyday structure of receiving an in-marrying stranger into a home, wanting to welcome without reserve and unable to cease being the one who welcomes, who therefore sets the terms. Every grammar of the rites is a particular, partial resolution of that impossible demand.

\subsection{The cross-cultural dyad: generating a ``third set of rites''}\label{sec:thirdrites}
When two people come from different grammars of rite, their diplomacy must coordinate two systems that do not map onto one another cleanly. The merely defensive response is translation---each learns enough of the other's grammar to avoid the gross solecism. But the \emb{generative} possibility is more than this: that the dyad fashions, between them, a \emb{third set of rites} belonging to neither family of origin but to themselves---a new, shared grammar in which both can stand and which neither inherited. This is the generative problem of the whole paper in its purest form, and the positive counterpart to \xref{sec:power}'s diagnosis. For \xref{sec:power} showed the default outcome to be asymmetric: not a third grammar but the victory of one, the marginalization of the other along the cultural-capital axis. The generative question is whether that default can be escaped.

Two failure modes must be named, because the third set of rites is easy to counterfeit. The first is \emb{dominance disguised as synthesis}: one side's grammar quietly furnishes the deep structure---the order of precedence, the real occasions, the language of the vows---while the other supplies only surface ornament, a borrowed dish, a token ceremony, a word of the other tongue. This looks like a third grammar and is the old asymmetry in fancy dress; the test is which grammar governs when the two conflict, and whose practice needs no explanation. The second is \emb{rootless pastiche}: a grammar assembled from fragments of both, observing the deep obligations of neither, so thin that it commands no grammaticality judgments and so generates no felt wrongness when violated---which is to say, not a grammar at all, since a grammar one can break without anyone wincing is a dead letter. A genuine third set of rites must do what \pinkref{§8.1} required of any grammar: combine, generate, and sustain judgments of fitness---and it must do so while drawing on two parent grammars without being captured by either.

Whether this is achievable, or whether the prior power asymmetry so structures the encounter that every apparent synthesis collapses into the first failure mode, is a question this paper \emph{poses and deliberately leaves open}. It is the hinge between the critical and the generative halves of the whole programme: if a third grammar can be made, then the domination \xref{sec:power} traced is not fate, and the rites can be remade to keep the care while loosening the rank (the hope of \pinkref{§7}'s coda); if it cannot, the honest conclusion is harder. The question is marked here as the site of a dedicated later paper, which a Sino-Japanese situation could serve, unnamed, as its working thought-experiment. \todo{develop the Sino-Japanese thought-experiment in the dedicated paper; keep unnamed here}

\goldrule
\footnotesize
\begin{longtable}{Y Y Y Y}
\caption{The rites, the sign, and the cross-cultural: literature}\label{tab:rite}\\
\toprule
\rowcolor{pageblush}
Source & Core contribution & Significance for the programme & Limitation / controversy \\
\midrule
\endfirsthead
\toprule
\rowcolor{pageblush}
Source & Core contribution & Significance for the programme & Limitation / controversy \\
\midrule
\endhead
Mauss, \emph{The Gift} (1925) & The three obligations; \emph{hau} & The micro-justice of the gift economy & Generalizing from archaic cases \\
Lévi-Strauss, \emph{Elementary Structures} (1949) & Exchange as structure; alliance & Marriage as exchange between groups & Structuralist universalism contested \\
Saussure; Peirce & Foundations of semiotics & The rites as a sign-system & Code-fixity overstated \\
Derrida, \emph{Of Hospitality} (1997) & The aporia of hospitality & The primordial ethics toward the other & Performative obscurity \\
Lévinas, \emph{Totality and Infinity} (1961) & The ethics of the face & The ethical face of the gaze & Asymmetry of the ethical relation \\
\bottomrule
\end{longtable}

% ============================================================
\section{The Symbolic (III): The Double Grammar of Law and Rite}\label{sec:law}

The external relational field is structured at once by the rites and by \emb{formal law}: marriage law, inheritance, the duty of maintenance, the legal nature of brideprice, the legal standing of in-laws, the law of divorce and of domestic violence. Law is the \emb{hard grammar} of the Symbolic; it coexists with the rites (the soft grammar), strains against them, and converts into them and they into it. Because law encodes and entrenches power asymmetry---in patrilocal default rules, in the conventions of surname transmission, in the default allocation of marital property---this chapter couples directly to \xref{sec:power}.

\subsection{Rite and law: the boundary of the informal and the formal}
Some diplomatic obligations are justiciable: maintenance, inheritance, the division of property. Others are upheld only by the rites and by social sanction: the return gift, filial deference, the attendance owed at a funeral. The boundary is neither fixed nor sharp, and the interesting cases are those of \emb{conversion}: a brideprice given as a rite becomes, when the engagement breaks, a debt the courts will hear---the gift crosses, under stress, from the soft grammar into the hard. To map which obligations sit on which side, and what pressures move an obligation across the boundary, is a research programme in itself.

\subsection{Legal pluralism}
Family ``quasi-law''---house rules, lineage regulations---coexists and conflicts with state law; and the cross-cultural dyad faces \emb{two jurisdictions at once} (the conflict-of-laws problems of a transnational marriage: which state's marriage law governs, which recognizes the union, how property and custody resolve across borders). The dyad here is not under one normative order but at the intersection of several, and its diplomacy includes the navigation of their conflict.

\subsection{Law as the hard coding of power}
Default rules are not neutral. The party who must act to depart from a default bears a cost the party favored by the default does not; and defaults in family law have historically been gendered---in residence, in name, in the presumption of property. Law thus hard-codes into enforceable form the asymmetries that \xref{sec:power} traced in the soft grammar. This connects to the author's prior work on sequential legal judgment (the NeSy-POMDP framework): the legal characterization of a diplomatic act is itself a sequential, partially observable inference, and the line of \xref{sec:bridge} reaches into the courtroom. \todo{supply precise jurisprudential entries; align with the author's existing legal-judgment work}

\goldrule
\footnotesize
\begin{longtable}{Y Y Y Y}
\caption{Jurisprudence: literature}\label{tab:juris}\\
\toprule
\rowcolor{pageblush}
Source & Core contribution & Significance for the programme & Limitation / controversy \\
\midrule
\endfirsthead
\toprule
\rowcolor{pageblush}
Source & Core contribution & Significance for the programme & Limitation / controversy \\
\midrule
\endhead
Legal pluralism (Griffiths 1986; Merry 1988) & Coexisting legal orders & The layering of rite / law / quasi-law & ``Law'' over-extended \\
Marriage law \& social reproduction & Law's shaping of the family & Law encodes power asymmetry & Jurisdictional variance \\
Brideprice as conditional gift (doctrine) & Rite converting to debt & A boundary case of rite/law & Doctrine differs by system \\
\bottomrule
\end{longtable}

% ============================================================
\section{The Symbolic (IV): The Creation and Circulation of Value---Political Economy}\label{sec:polecon}

Political economy enters this paper in a register it is rarely given: a \emb{positive, constructive} one, not merely critical. The reflex, in a critical tradition, is to treat the political economy of intimacy as a story of extraction alone---brideprice as commodification, family labor as exploitation. That story is true and \xref{sec:justice} will tell it; but to begin there is to miss what generative justice insists upon, which is that value is first \emph{produced} and can \emph{circulate} before it is, or instead of being, extracted. Diplomatic commerce creates value---relational capital, trust, recognition, affect, social network---and makes it circulate through chains of reciprocity, the flow of the gift, the spread of reputation; against which, and only then, stands value extracted in one direction. This is the positive battlefield of Eglash's generative justice, and the most natural redemption of the Value Foam framework in the diplomatic field: circulation is value flowing along the field and returning to those who generate it; extraction, alienation, is holonomy---value returning to its origin unconserved after a circuit of an asymmetric field. The chapter takes the two in order: first what is made, then how it circulates, then how it leaks.

\subsection{What value diplomacy creates}
The first claim is the easiest to overlook precisely because it is positive: that the exchanges of diplomacy \emph{generate} goods that did not pre-exist them. These goods are real though not monetary. There is relational capital---the standing accumulated by reliable reciprocity, the credit one can draw on because one has given before. There is, in Bourdieu's terms, social capital (the network one can mobilize) and symbolic capital (the honor, the recognized standing, one is accorded). And there is something Bourdieu's vocabulary, with its faint odor of strategy, undersells: the trust and affection themselves, which are not merely instrumental resources but goods internal to the relation, valuable in being had and not only in being usable.

That these are genuinely \emph{produced}, and not merely transferred, is the load-bearing claim, and it distinguishes a generative account from a zero-sum one. A wedding does not redistribute a fixed stock of family goodwill; it can \emph{create} goodwill that existed in neither family before, a new relation between two lineages that is a positive addition to the world's stock of relation. A well-conducted season of diplomacy leaves more trust in being than it found. This is why the political economy of intimacy is not exhausted by its critique: there is a positive product, and the normative question---the generative-justice question---is whether that product \emph{circulates} to those who made it or is \emph{extracted} from them. One cannot even pose the question of just circulation until one has seen that there is something produced to circulate. The critical tradition's haste to extraction skips this step and thereby loses the standard against which extraction is a wrong: extraction is unjust precisely because it diverts a value that could have circulated.

\subsection{The mechanisms of circulation}
Value, once produced, circulates by definite mechanisms, and naming them lets one say with some precision what healthy circulation is. It moves through \emb{chains of reciprocity}: the gift answered by a return gift answered in turn, each link both discharging a debt and opening a new one, so that the chain is self-perpetuating and binds its parties over time. It moves through the \emb{flow of the gift across the network}: a thing received and passed on rather than retained, the Maussian circulation in which the gift's value lies in its movement, and to hoard it would be to kill it. It moves through the \emb{spread of reputation}: standing earned in one quarter travels to another, so that the value of having given well is realized far from the giving, in the credit one is extended by strangers who have only heard.

Where these mechanisms run freely, value is retained \emph{within the network and by its producers}---it returns, by some path, to those whose giving generated it---and the network is, in the generative sense, just. This is the positive ideal against which the next subsection's failure is measured: not equality of holdings but \emph{circulation}, value kept in motion among those who make it rather than pooled at a point of accumulation. Generative justice is, in this register, a fluid-dynamic rather than a distributive notion: the question is not who holds how much at an instant but whether the value flows, and to whom it returns over a circuit.

\subsection{Extraction and alienation}
Circulation has its shadow, and the shadow is the chapter's critical payload---reached now, after the positive ground, so that it lands as the corruption of something rather than as the whole story. \emb{Emotional labor}---the work of maintaining relations, remembering occasions, anticipating needs, smoothing slights, holding the family's affective fabric together---is, characteristically, \emb{gendered}: borne disproportionately by women and by the in-marrying party, and frequently uncompensated by any return. Here the value produced does not circulate back to its producer; it \emb{leaks} to the superior party, who enjoys the relational goods---the smoothed gathering, the remembered birthday, the managed tension---that another's labor produced, while the producer's standing does not rise by it and her effort goes unmarked because it is presumed natural to her position.

This is the diplomatic-field instance of the generative-justice diagnosis, and the Value Foam formalism states it exactly: the leakage is the \emb{non-conservation of value over a circuit of the asymmetric field}. Recall \xref{sec:power}'s result that on a tilted field the connection is asymmetric; the consequence here is that value carried around the loop does not return to its origin. The holonomy is non-trivial, and what the weaker party puts in does not come back to her---it has been parallel-transported, by the field's own asymmetry, to the stronger. Alienation, in this register, has a precise meaning, neither merely Marxian nor merely metaphorical: it is value \emph{generated by one and accumulated by another} by the geometry of an unequal field, with no theft at any single step, the leakage distributed invisibly across a circuit each of whose links looked fair. This is why the wrong is so hard to name from inside: there is no moment of taking to point to, only a loop that does not close. The formal demonstration---the connection made explicit, the conditions for non-trivial holonomy, the magnitude of the leak---is the work of the dedicated Value Foam paper. \todo{confirm precise Eglash entry; align Value Foam notation with the series' prior formulation}

\goldrule
\footnotesize
\begin{longtable}{Y Y Y Y}
\caption{Political economy and the circulation of value: literature}\label{tab:polecon}\\
\toprule
\rowcolor{pageblush}
Source & Core contribution & Significance for the programme & Limitation / controversy \\
\midrule
\endfirsthead
\rowcolor{pageblush}
Source & Core contribution & Significance for the programme & Limitation / controversy \\
\midrule
\endhead
Marx, theory of value; \emph{1844 Manuscripts} & Use/exchange value; alienation & The fourth, relational value & Labor theory contested \\
Eglash, generative justice & Circulation vs.\ extraction & The chapter's main line & Precise entry to be confirmed \\
Bourdieu, ``The Forms of Capital'' (1986) & Social / symbolic capital & The convertibility of relational capital & Economism charge \\
Social reproduction theory (e.g.\ Federici) & Unwaged care labor & The gendered structure of extraction & Scope of ``reproduction'' \\
\bottomrule
\end{longtable}

% ============================================================
\section{The Mechanism Bridge: The Cognition and Decision of Relational Value}\label{sec:bridge}

This chapter cuts across the Symbolic and the phenomenal, supplying a \emb{computational-neural chain}: the relational state (a latent variable) $\rightarrow$ the neuro-cognitive system performs approximate Bayesian \emb{state inference} upon it (the prefrontal cortex, active inference) $\rightarrow$ under partial observability it makes a \emb{sequential decision} (POMDP, games) $\rightarrow$ the decision is the diplomatic act (the rites of the phenomenal field). This chain threads the structure of the Symbolic, the cognition of value, and the acts of the phenomenal field onto a single line---and it is one of the paper's distinctive contributions, the thing a philosopher without the author's computational-neuroscientific background could not readily supply.

\subsection{The cognition of value: higher-order mentalizing}
The chapter's chain begins in cognition, because before the dyad can act in the field it must \emph{read} the field, and reading it means inferring how others appraise the relation. This is a mentalizing problem, and higher-order coupling makes it a problem of \emb{nested} mentalizing of ascending order. It is not enough to model what one's partner thinks; one must model what one's mother thinks of one's partner, and---a further order---what one's partner imagines one's mother to think of them, and, in the heat of a difficult dinner, what one's mother is reading in the partner's face as the partner tries to manage exactly that imagining. This recursion (third-, fourth-order, and on) is the precise cognitive content of the slogan ``from second-order to higher-order coupling'': the order of the coupling is, cognitively, the depth of the nesting the situation demands.

Two features of this nesting are worth drawing out, because they explain phenomena the slogan alone does not. \emc{First, the load is combinatorial, not linear.} Adding one third party does not add one mental model; it adds a tree of models---what they think of us, what we think they think, what our partner thinks they think of us---whose branching is why the entry of even a single in-law can feel disproportionately, almost vertiginously, taxing in a way that the arithmetic of ``one more person'' fails to predict. The felt difficulty of the first family dinner is not innumeracy; it is the genuine combinatorial cost of the tree. \emc{Second, the recursion does not terminate cleanly.} In principle the nesting regresses without end (what she thinks I think she thinks\ldots); in practice it is truncated, and the question of \emph{where} a culture or a person truncates it is itself diagnostic. Some fields tolerate stopping at the second order (``I needn't guess what she imagines I feel; I will simply ask''); others demand the fourth, and punish a truncation read as obtuseness. The depth of mandatory mentalizing is, in this sense, a cultural and a power-laden variable, not a fixed cognitive ceiling---and the party with less standing is typically the one required to mentalize deeper, since they bear the burden of anticipating a reading they cannot afford to get wrong. This extends the joint-attentional account of \pinkref{Paper~V} in a specific way: the shared attentional field of the dyad must now hold not only the world the two attend to together but a recursive representation of how third parties attend to \emph{them}---attention turned reflexive, and turned upward along the gradient of power.

\subsection{Neuroscience: the implementation of state inference (cautiously)}
\emc{The stance here must be stated before the content, because the content invites misreading: what follows is offered as one possible \emph{implementation} of the mechanism, not as a reduction of relational value to neural activity.} The distinction is the familiar one between levels of description, and the chapter's whole defensibility rests on keeping it. With that fixed, the candidate substrates are well attested. The medial prefrontal cortex (mPFC) is implicated in mentalizing and in the representation of self- and other-related value; the temporo-parietal junction (TPJ) in the attribution of beliefs; the orbitofrontal and ventromedial prefrontal cortex (OFC/vmPFC) in the encoding of value and social reward; the dorsolateral prefrontal cortex (dlPFC) in norm-compliance and fairness-related decision, as in the neural studies of the ultimatum game. Above these particulars sits the unifying frame: the relational state is a \emb{latent variable} upon which the subject performs approximate Bayesian inference---active inference, in Friston's free-energy formulation, in which perception and action alike work to minimize expected surprise. The relational state is never directly observed; it is inferred from signs (a tone, a gift's thickness, an absence), and the inference is revised as new signs arrive. This is the same formal posture---latent state, noisy observation, sequential update---that governs the author's computational-neuroscientific work on EEG and grammar-based modeling at NCNP, and the recurrence of that posture across such different domains is part of what recommends it.

Two guards are needed, one against deflation, one against inflation. \emc{Against reduction:} that a process is implemented in cortex does not make the relational value \emph{nothing but} cortical activity, any more than a sentence's being realized in ink makes its meaning nothing but ink, or a proof's running on silicon makes the theorem a fact about transistors. The implementational level answers ``how is it physically realized,'' not ``what is it''; the value of a recognition, the wrong of a slight, belong to the relational and normative levels, which the neural level realizes without exhausting. To forget this is to commit the error the chapter's own framing (\xref{sec:frame}, \pinkref{§2}) forbids: letting one stratum usurp the others. \emc{Against mere metaphor, in the other direction:} it would be too cheap to invoke ``Bayesian inference'' and ``prefrontal cortex'' as decorative gestures toward rigor. The claim earns its keep only if the formal posture does work the looser vocabulary cannot---and it does, in three places the later computational paper will develop: it predicts that diplomatic missteps cluster where the \emph{signs are noisiest} (ambiguous gifts, silences readable two ways), not where the stakes are highest; it explains the felt cost of a new in-law as the cost of inference under a suddenly enlarged latent space; and it makes ``tact'' a precise notion---the policy that minimizes expected surprise for \emph{all} parties' inferences at once, rather than only one's own. A frame that yields predictions, an explanation, and a definition is not a metaphor. It is the implementational layer of the chain whose next link is decision.

\subsection{Decision theory: from state inference to the diplomatic act}
Inference issues in action. Once the relational state is inferred---however provisionally, under whatever noise---the subject must \emb{decide}: what to give, whether to attend, how to answer an affront, when to let a slight pass. The natural formalism is sequential decision under partial observability, a POMDP, the relational state being precisely partially observable: one acts not on the state but on a \emph{belief} about the state, updated as signs arrive, and a good diplomatic policy is one that performs well across the spread of states one's belief allows. This is not an analogy borrowed from elsewhere; it is the same structure the author's NeSy-POMDP work brings to sequential legal judgment, and its recurrence here is the chain's link from cognition into act.

The value of the game-theoretic vocabulary is that different game-forms isolate different things the looser word ``interaction'' blurs together, and each names a real diplomatic situation. \emc{Coordination games} model the case where the parties' interests align but the equilibrium is underdetermined: two families who both want one wedding must converge on \emph{a} form, and which form matters less than that they meet---the pure problem of the cross-cultural double ceremony (\xref{sec:network}), where the difficulty is not conflict but the absence of a default. \emc{Signaling games} model the gift as a \emb{costly signal}: a gift communicates commitment precisely because it costs, so that what is transmitted is not the object but the credibility the cost underwrites---which is why a cheap gift given easily and a dear one given at sacrifice say different things though both are ``gifts,'' and why a gift that costs the giver nothing fails to signal at all. \emc{Repeated play} models reciprocity: the back-and-forth of return gifts is a culturally elaborated reciprocal strategy, in which today's apparent loss is tomorrow's claim, and defection (the unreturned gift, the unanswered visit) is punished by the slow withdrawal of the relation. Each form picks out a structure the others miss; together they are why ``diplomacy'' is not one game but a field of them, played at once.

One correction to the standard apparatus is required, and it is substantive. The agent of diplomacy does \emph{not} maximize narrow self-interest; decision proceeds under \emb{social preferences}---inequity aversion, a taste for reciprocity, an aversion to having wronged---so that a diplomatically rational agent will pay, in money or face or effort, to right an imbalance that a self-interested agent would happily leave standing. This is not a departure from rigor but a correction of the utility function toward the one the phenomena actually display; a model that predicted indifference to a debt owed would simply be wrong about what people do. With that correction, the chain closes: the structure of the Symbolic conditions what states are possible and how they are read, inference under that structure yields a belief, decision under social preferences turns the belief into a policy, and the policy's output \emph{is} the diplomatic act---the gift given, the seat taken, the visit made or withheld---that the phenomenal chapters (\xref{sec:triad}, \xref{sec:network}) now take up as their material. The bridge is built; we cross it into the phenomenal field.

\goldrule
\footnotesize
\begin{longtable}{Y Y Y Y}
\caption{Cognition, neuroscience, and decision: literature}\label{tab:neuro}\\
\toprule
\rowcolor{pageblush}
Source & Core contribution & Significance for the programme & Limitation / controversy \\
\midrule
\endfirsthead
\toprule
\rowcolor{pageblush}
Source & Core contribution & Significance for the programme & Limitation / controversy \\
\midrule
\endhead
Friston, active inference / free energy & The Bayesian brain & A unifying frame for state inference & Falsifiability debated \\
mPFC / TPJ mentalizing studies & Neural bases of social cognition & Implementation of higher-order mentalizing & Reverse-inference, replicability \\
OFC / vmPFC value coding & Value and social reward & Neural representation of relational value & Reductionism risk \\
POMDP / sequential decision (Sutton \& Barto) & Decision under partial observability & Formal model of the diplomatic act & Computational tractability \\
Signaling (Spence); reciprocity (Fehr) & Costly signals; social preferences & The gift as signal; reciprocity & Idealized assumptions \\
\bottomrule
\end{longtable}

% ============================================================
\section{The Phenomenal (I): The Triad---The First Other Beyond the Dyad}\label{sec:triad}

We cross now from the strata that condition diplomacy into the field where it is enacted, and we begin at the smallest possible scale of the higher-order: the dyad plus one third party. This is the minimal higher-order coupling, and it is chosen as the entry to the phenomenal not for simplicity but because it is the threshold at which the genuinely new appears. There is a deliberate resonance with the three-body problem of mechanics: just as three gravitating bodies admit no closed solution though two do, the triad is not a manageable extension of the dyad but a qualitative break, the point at which the relational system loses the closed-form tractability it had at second order. One third party is enough to change the kind of problem.

The entry of that third party \emb{activates three things at once}: the triangular structure (psychoanalysis), the ascent of nested mentalizing (the cognition of \xref{sec:bridge}, \pinkref{Paper~V}), and the dyad's first experience of \emb{being watched as a whole} (the gaze of the Other). The claim that they emerge \emph{together}, symmetrically, rather than one giving rise to the next, is the methodological point of \pinkref{§2} made concrete: none of the three is the master phenomenon from which the others derive. One could not say the triangle causes the gaze, or the mentalizing grounds the triangle; they are three faces of one event, the entry of the third, and to privilege any as primary would be the hub-error the whole paper refuses. The triad is thus the first place the paper's decentered method is not merely asserted but exhibited---three co-equal phenomena, organized by their coupling at the single node ``a third has entered,'' and by nothing more.

\subsection{The psychoanalytic structure of the triangle}
The triangle here is not, or not only, the sexual triangle of rivalry for a beloved. The structures ``mother--son--daughter-in-law'' and ``father--daughter--suitor'' are triangles in which a member of the dyad is positioned between the partner and a prior, primary bond---and in which jealousy, exclusion, and the strain of mediation play out along lines the participants did not draw. The reason a third party necessarily produces a triangle, rather than merely a third relationship, is structural: where there were three pairwise bonds there is now a closed figure, and a closed figure has an inside and an outside, a position that is "between" and positions that are "joined against." Three points cannot all be adjacent without one of them being, on any given question, the odd one out.

What makes this more than family-systems folklore is its connection to the Real of \xref{sec:real}. The entering partner is liable to be cast---without anyone choosing it, often against everyone's conscious wish---as the one who has "taken" the son, or as the surface onto which a mother's own unsymbolized grief is transferred. The partner becomes the third term whose arrival destabilizes a dyad that had its own long, unspoken accounts, and is made to carry a charge that was there before them and is not about them. This is the elementary form in which the dyad's boundary is tested from without: not by argument over stated matters, but by the silent assignment of a position in a triangle whose tensions predate the marriage. The diplomacy of the triangle is therefore rarely the diplomacy of negotiation; it is the harder work of refusing a position one is being maneuvered into without ever being addressed.

\subsection{The ascent of nested mentalizing}
The triad demands, in lived form, the recursion that \xref{sec:bridge} described abstractly: "I imagine how my partner imagines my mother's view of them." In the dyad this recursion was optional and shallow; in the triad it becomes obligatory and deep, because action in the presence of the third is governed by it. What one says, what one leaves unsaid, whom one defends and how visibly, when one reaches for a partner's hand under a parent's eye---each is a move chosen against a model of how it will be read by two parties at once, each of whom is reading the other reading it. The combinatorial cost named in \xref{sec:bridge} is here not a theoretical point but the felt texture of the first difficult dinner: the exhaustion is real, and it is the exhaustion of holding a branching tree of models live under time pressure.

Two consequences follow that the dyadic account could not yield. First, the recursion creates the possibility of a distinctively triadic failure: \emb{miscoordination through over-mentalizing}, where each party, modeling the others so intently, acts on a predicted reaction that the prediction itself provokes---the partner who, anticipating the mother's disapproval, stiffens in a way that produces the disapproval anticipated. The triangle can manufacture the very tensions its members are straining to avoid, a feedback the dyad's simpler structure does not permit. Second, the burden of this mentalizing is, as \xref{sec:power} warned, unequally distributed: the lower-standing party must mentalize deeper, since they bear the cost of misreading a party they cannot afford to offend, while the higher party may decline the labor altogether and simply be read. To mentalize hard is itself a mark of the weaker position; the one who need not guess what others think of them is the one with standing.

\subsection{Being watched as a whole}
In the dyad each is watched by the other; in the triad the dyad is, for the first time, watched \emb{as a unit}. "Are they well matched?"---the question is not about either partner but about the couple, which is thereby constituted as a single object of judgment, no longer two individuals appraised separately but "the couple," held up against an image of what a fitting couple should be. This is a genuine ontological addition, not merely a new opinion about existing entities: the gaze of the third \emph{brings into being} the couple-as-object, a thing that did not exist while only the two regarded each other.

To see why this is more than metaphor requires the phenomenology of intersubjectivity. On Husserl's account \citep{husserl1931cm}, the other is given to me not as an inferred object but as an \emph{alter ego}, a second center of experience before whom I too become appresented as a body-for-others; the third party's entry is thus the entry of a new constituting consciousness in whose experience the couple takes form. Schutz's \emph{we-relation} \citep{schutz1932} sharpens the point for this paper specifically: the dyad is a lived ``we,'' built in the face-to-face mutual tuning-in of two streams of consciousness, and the triad forces that we-relation to either expand to include the third or to harden, under the third's gaze, into a ``they'' observed from outside. The couple-as-object is the we seen from a position no member of the we occupies---which is why it can feel like a stranger's possession of something that was theirs.

That this objectification is two-edged is the crux. It is a source of \emb{recognition}: to be seen as a couple is to have the union acknowledged into existence, registered in the symbolic order of others, made real in a way the dyad cannot make itself real from inside. A union no one sees is, in a sense the lovers may find unbearable, not yet fully actual. And it is a source of \emb{alienation}: to be seen as a couple is also to be reduced to the image others hold, flattened into "that couple," fixed as an object whose meaning is now partly out of the partners' hands. The paper locates itself, here, between two accounts of the gaze it declines to choose between in the abstract. Sartre's look objectifies \citep{sartre1943}: it fixes the self as an object for another, induces shame, steals one's freedom by making one a thing in someone else's world. Lévinas's face does the opposite \citep{levinas1961}: the other's regard is not a threat but an ethical summons, a claim that calls one to responsibility before it calls one to defense. The thesis of this section is that the gaze of the Other in diplomacy is \emph{neither purely one nor the other}, and that which it becomes is not fixed by the gaze itself but by the \emb{structure within which it falls} (\xref{sec:power}). Here the \emph{critical} phenomenology of the gaze does decisive work: Ahmed's analysis of \emph{orientation} \citep{ahmed2006queer} shows that bodies are not appraised from nowhere but along inherited lines---some couples are ``in line'' with what the room expects and fall easily into the field of approving recognition, others are ``oblique,'' out of line, and meet a gaze that objectifies before it ever assesses the particular two. The appraising look of an elder who holds the power of acceptance, down the gradient, tends toward the Sartrean: it objectifies because it judges from a height the judged cannot return. The look of a friend who meets the couple as equals can be the Lévinasian summons, recognition without subjection. The gaze is not good or bad in itself; it takes the moral character of the field it crosses. This is why the same wedding can feel, to a couple, at once like being blessed and like being inspected---because it is both, from different positions around the same room.

\begin{casebox}{The in-law triangle}
A partner is drawn into the unfinished emotional business of the other's family of origin. A mother and her son have a long dyad with its own unspoken accounts; the entering partner becomes the third term, and finds themselves cast as the one who has ``taken'' the son, or the one onto whom the mother's own old griefs are transferred. The partner is mediator and target at once, pulled between loyalty to the dyad and the pressure of the prior bond. The case activates: the triangle, transference (\xref{sec:real}, the traumatic real of the family), and the axes of power.
\end{casebox}

\begin{casebox}{The appraising look at the wedding}
At the wedding the couple is appraised as a whole---``are they well matched?''---by a gathering whose gaze measures, judges, and silently records. The two feel themselves become a single object, their union held up against an image of what a fitting union should look like. The case activates: the gaze, the imaginary register of the fitting image, and value (the couple appraised as if it had a worth that could be read off).
\end{casebox}

\goldrule
\footnotesize
\begin{longtable}{Y Y Y Y}
\caption{The triad: literature}\label{tab:triad}\\
\toprule
\rowcolor{pageblush}
Source & Core contribution & Significance for the programme & Limitation / controversy \\
\midrule
\endfirsthead
\toprule
\rowcolor{pageblush}
Source & Core contribution & Significance for the programme & Limitation / controversy \\
\midrule
\endhead
Lacan, the gaze and the big Other & The subject constituted in the Other's look & The mechanism of being-watched-as-a-whole & Textual obscurity \\
Sartre, \emph{Being and Nothingness} (the look) & Objectification and shame & The objectifying face of the gaze & Pessimism of the for-others \\
Family systems theory (Bowen, triangles) & The triangle as a stabilizing device & The dynamics of the in-law triangle & Empirical basis contested \\
\bottomrule
\end{longtable}

% ============================================================
\section{The Phenomenal (II): Network and Field---The Enactment of Gift and Rite}\label{sec:network}

From the triad the ascent continues to the many-bodied: to network and field, where not one third party but a whole structured surround---two families of origin, a circle of friends, a network of colleagues---bears upon the dyad at once. Here the phenomena the paper has named throughout finally take the stage as acts: the gift, brideprice and dowry, forms of address, seating and toasts, the politics of absence, the cross-cultural double ceremony. The governing instruction for reading them is the one \xref{sec:method} set and the preceding chapters earned: each is to be read as the \emb{enactment, in the phenomenal field, of the power and grammar of the Symbolic} (\xref{sec:power}--\xref{sec:polecon}). Because the Symbolic has been laid out first, these acts arrive not as isolated anecdotes---curiosities of custom to be catalogued---but as \emph{utterances in a grammar}, moves on a tilted field, surface expressions of a structure the reader already holds. An anecdote illustrates; an enactment manifests. The difference is the whole justification of the paper's order.

This is also where the formal vocabulary of optimization and dynamical systems enters, no longer as a framework listed (\xref{sec:frame}) but as a reading applied at the level of behavior, which \pinkref{§13.5} develops once the acts are before us.

\subsection{The phenomenology of gift and brideprice}
The gift is where the four-fold value of \xref{sec:rite} becomes visible action. Its \emb{thickness} is measured, compared against precedent, entered silently into the receiving family's running ledger---and the giver, knowing it will be weighed, chooses under that knowledge, so that the gift is already a move anticipating its own appraisal. Its \emb{timing} speaks: a return gift too prompt reads as the discharge of a burden, an eagerness to be quit of obligation that insults by refusing the ongoing bond a debt sustains; too slow reads as neglect, the obligation forgotten or disdained. The interval itself is a message, independent of the gift's content, because in a grammar of reciprocity the tempo is part of the syntax. And the gift's \emb{figure}---the sum on a red envelope---speaks a semantics of auspicious and inauspicious numbers, in which the amount is read not as quantity but as sign, a wrong number souring a right sum.

The decisive analytic point is that these are not \emph{decorations on} an underlying economic transfer, a layer of meaning painted over a movement of money. They \emph{are} the diplomatic act; the transfer is their vehicle, not their substance. Brideprice makes this sharpest, and shows why the same act sits on the boundary \xref{sec:justice} will probe: as recognition it is a costly, public sign of the worth of a union and of the family that raised the bride, the gift's relational value (\xref{sec:rite}) made ceremonial; as commodification it slides toward the pricing of a person, the conversion of relational value into a sum, which touches the inalienable (\xref{sec:real}) and wounds. Nothing in the transfer itself decides which it is; the same payment is honor or purchase depending on the grammar read into it and the structure (\xref{sec:power}) it moves within. The gift is thus the node where every stratum of the paper meets in a single act---which is exactly why \pinkref{§2} refused to let it become the hub: it is richly connected, but its richness is the meeting of the other dimensions in it, not a source from which they flow.

\subsection{Address, seating, toasts: the power-coding of language and space}
A class of diplomatic acts carries almost no object and yet enacts the hierarchy of \xref{sec:power} more exactly than any gift, precisely because it passes as "mere form." How one \emb{addresses} a partner's parents is a positioning: the kin term assigns a place in the order and to utter it is to accept that place, so that the choice of word is the acceptance of a rank---which is why it is so fraught where two grammars meet (the case below). \emb{Seating} encodes rank in space: the head and foot of a table, the proximity to the honored guest, are a diagram of standing that everyone reads and no one states. \emb{The order of toasts} encodes rank in time: who is honored first, who must rise to whom, in what sequence, lays out the hierarchy as a temporal sequence performed before all. These wordless or near-wordless rites are the purest enactment of structure, and their power lies exactly in their passing as empty form: because "it is only where one sits," the hierarchy it diagrams is naturalized, placed beyond contest, made to seem a fact of arrangement rather than a claim of rank. The acts that look least like assertions of power are often its most efficient vehicles, since what is not said as a claim cannot be answered as one.

\subsection{The politics of absence}
Presence is the medium of recognition; absence is therefore its most charged withdrawal. Who is not invited, who fails to appear at a wedding or a funeral, is a diplomatic utterance, and frequently the loudest one in the field---louder than any said thing, because it cannot be softened in the saying. To withhold presence is to withhold recognition; a calculated absence is a declaration that needs no words and admits no easy reply, for to protest an absence is to confess one wanted the presence. The empty seat at the wedding, the elder who does not come, the child who stays away from the rite, signals a rupture or a refusal that all present read instantly and no one may comfortably name. Presence and absence form the binary in which the field's deepest alignments and breaks are signaled---a one-bit message of enormous force, precisely because its very crudeness makes it unanswerable. This is the gift's logic inverted: where the gift signals by what is given, absence signals by what is withheld, and withholding, having no content to dispute, is the harder utterance to counter.

\subsection{The cross-cultural double ceremony}
Two ceremonies, two grammars, coordinated into one marriage: this is the phenomenal realization of the "third set of rites" posed in \xref{sec:thirdrites}, the place where that abstract question takes visible form. And the visible details are exactly the test \pinkref{§8.4} specified for distinguishing genuine synthesis from its two counterfeits. Which ceremony is treated as the "real" one and which as the courtesy? Whose kin set the terms, the date, the seating? In which language are the vows spoken, and which side must follow in translation? These are not logistical trivia; they are the diagnostic surface on which one reads whether a true third grammar is being generated or whether one grammar furnishes the deep structure while the other supplies ornament (the first failure mode), or whether the two are merely juxtaposed without either's deep obligations honored (the second). The double ceremony is where the generative hope of the whole paper is either realized or quietly defeated, and the defeat, when it comes, is legible not in any quarrel but in whose practice needed no explaining.

\subsection{The behavioral reading in optimization and dynamics}
With the acts before us, the formal vocabulary of \xref{sec:frame} can be applied at the level of behavior, and it earns its place by making precise two things the prose has so far said loosely. The first is what "decorum" \emph{is}. The diplomat faces conflicting objectives at once---internal accord within the dyad, the expectations of each family of origin, the ambient social norms, the dyad's own values---and these do not jointly maximize; there is no single act that best satisfies all, only trade-offs among them. In the language of multi-objective optimization, there is no global optimum, only a \emb{Pareto frontier} of acts each of which sacrifices on one objective to gain on another. "Decorum," on this reading, is not a code one obeys but the \emb{feasible set}: the region of acts that violate no hard constraint and lie on or near the frontier, within which the diplomat must still choose, trading a little of the family's expectation against a little of the couple's own values. This dissolves a confusion in the ordinary notion of the "proper": there is rarely one proper act, because the objectives conflict; there is a set of defensible ones, and skill is the navigation among them.

The second is the field's \emb{dynamics}. The relational field is not static; it evolves, and each gift and slight is a perturbation to a system with its own tendencies. Debt and reciprocity are \emb{feedback loops}---a gift creates a pressure toward return, a slight a pressure toward retaliation or withdrawal---and a field of such loops has stable basins and unstable ridges. The art of diplomacy, in these terms, is to keep the system in a stable basin without letting it rigidify: a relation can fail by collapsing into open conflict (the loops running to rupture) but equally by \emph{rigidifying} into pure form, a hollow correctness in which every rite is observed and no feeling moves, the dead grammar of \pinkref{§8.4.} Stability is not the goal; \emph{living} stability is---a basin deep enough to absorb the ordinary perturbations of slight and misunderstanding, shallow enough that feeling still flows. A worked optimization, with the objectives specified and the frontier drawn, belongs to the dedicated computational paper; here the reading is qualitative, its purpose to show that "decorum" and "tact" are not vague honorifics but names for solutions to a structured problem. \todo{worked toy optimization deferred to the computational sub-paper}

\begin{casebox}{The weighing of the first-visit gift}
On the first formal visit to a partner's family, the gift brought is measured on arrival---its thickness noted, compared with precedent, entered into the family's running ledger. The giver knows it will be weighed and chooses accordingly; the weighing is done, often, without a word. The case activates: value, the gaze, power (the right to measure), and the semantics of the sign.
\end{casebox}

\begin{casebox}{The problem of address}
A partner must decide how to address the other's parents---a choice that is no mere label but the acceptance of a position in the kin order, and that may differ across the two families' grammars, so that the term proper in one is presumptuous or cold in the other. The case activates: the symbolic grammar of the rites, power (the position the term concedes), and cultural difference.
\end{casebox}

\begin{casebox}{The taboo misgift}
A gift chosen with good intent---a clock, an umbrella, white flowers---carries, in the recipient's grammar, an association of death, separation, or mourning, and the giving becomes a hermeneutic incident: the act read as the opposite of what was meant. The case activates: the semantics of the sign, the right of interpretation, and the wound of misreading (\pinkref{Paper~VI}).
\end{casebox}

\begin{casebox}{Absence as declaration}
An elder's absence from a wedding, or a junior's from a memorial rite, is read by the whole field as a statement---of disapproval, of rupture, of a refusal to recognize. No accusation is made; the empty seat makes it. The case activates: presence/absence, power, and justice (the withholding of recognition).
\end{casebox}

\goldrule
\footnotesize
\begin{longtable}{Y Y Y Y}
\caption{Gift phenomena and optimization/dynamics: literature}\label{tab:phenom}\\
\toprule
\rowcolor{pageblush}
Source & Core contribution & Significance for the programme & Limitation / controversy \\
\midrule
\endfirsthead
\rowcolor{pageblush}
Source & Core contribution & Significance for the programme & Limitation / controversy \\
\midrule
\endhead
Gift-economy anthropology (Mauss; Sahlins) & The phenomenon of reciprocity & The gift at the level of behavior & Generalization across cultures \\
Multi-objective optimization & The Pareto frontier & A formal reading of ``decorum'' & Functionalist risk \\
Dynamical systems / attractors & Stability vs.\ rigidity & The evolution of the relational field & Metaphor vs.\ model \\
\bottomrule
\end{longtable}

% ============================================================
\section{The Tension of Coupling: Value, Debt, and Justice}\label{sec:justice}

This chapter is the one the whole stratal order was built to make possible. Justice could not be treated first, as a flat theory of fair exchange, because the diplomatic field is not flat; it had to wait until the Real had marked what cannot be priced (\xref{sec:real}), the Symbolic had revealed the gradient of power, the grammar, and the machinery of value (\xref{sec:power}--\xref{sec:polecon}), the bridge had shown how value is inferred and decided (\xref{sec:bridge}), and the phenomenal had displayed the acts in which all of this is transacted (\xref{sec:triad}--\xref{sec:network}). Only now can the treatment of justice \emb{integrate the three registers} rather than abstract from them. The integrating thesis is this: every diplomatic exchange generates asymmetry and debt, and upon the unequal field of \xref{sec:power} these must be \emb{re-read}. What the flat liberal picture sees as the fair give-and-take of equals is, on a tilted field, something else---and the chapter's task is to show that the three classical registers of justice (distributive, recognitive, hermeneutical) do not merely each apply but \emph{knot together}, tightest where the field is most unequal.

\subsection{The temporality of debt and reciprocity}
The gift obliges a return; this is the micro-justice of the gift economy \citep{mauss1925gift,graeber2011debt}, the symmetry that keeps a relation alive across time, since an unreturned gift would end the exchange and a fully settled account would end the bond. So far reciprocity is just, and even generative: the open debt is what keeps the relation in motion (\xref{sec:polecon}). But the obligation has a \emb{temporal structure that can be weaponized}, and seeing how requires the asymmetric field. Between equals, a large gift is answered, in time, by a comparable one, and the parties remain peers across the cycle. Where the parties are unequal, a gift \emph{too large to be returned} does not bind them as equals; it \emb{subordinates} the receiver, who now owes what cannot be repaid and must pay, in the only currency left, by deference. The grand gift, the unrepayable favor, the education paid for, the house provided---these can be pure generosity, and they can be a technique of domination wearing generosity's face, and the recipient often cannot tell which until the deference is called in.

The criterion that separates the two is precisely \emph{reparability}. Reciprocity is just only between parties able, in principle, to reciprocate; the symmetry of the gift economy presupposes a rough equality of capacity to give in return. Where that capacity is absent---because the gift is too large, or the receiver too poor, or the positions too unequal---the \emph{form} of reciprocity persists while its justice drains out, and what remains is a relation of dependence dressed in the etiquette of exchange. This is the temporal mechanism behind \xref{sec:power}'s economic axis and \xref{sec:polecon}'s leakage seen from the side of justice: the unrepayable gift is how value, given once, converts into standing held indefinitely. The just response to such a gift is not refusal (which insults) nor acceptance-as-subordination, but the difficult diplomatic work of converting it back into a relation among equals---which the asymmetric field makes hard, and sometimes forecloses.

\subsection{Recognition versus commodification}
Brideprice is the sharpest case because the same act faces two ways. As \emb{recognition} it is a public, costly acknowledgment of the worth of a union and of the family that raised the bride---value made ceremonial, the relational value of \xref{sec:rite} given a form the whole field can witness. As \emb{commodification} it is the pricing of a person, the conversion of a daughter into a sum, the reduction of a who to a how-much. The unsettling fact is that the two are not always distinguishable \emph{by the act alone}: the same transfer, the same amount, can be felt as honor or as purchase depending on the structure it moves within (\xref{sec:power}) and the meaning the parties read into it (\xref{sec:rite}). This is why the question cannot be settled by inspecting the custom in the abstract, as critics and defenders of brideprice alike attempt; it is settled, if at all, in the particular field and the particular reading.

At its limit commodification touches the Real (\xref{sec:real}): the wound of being priced is the wound of having one's relational, inalienable worth converted \emph{without remainder} into exchange-value, the category error of \pinkref{§6} felt as injury. And here the positive pole, recognition, carries its own danger, which the asymmetric field forces into view and which cannot be waved away. Honneth's theory \citep{honneth1992} supplies the ideal of recognition as the good opposed to disrespect; but recognition can itself be \emb{ideological}---it can "recognize" a person precisely \emph{into} a subordinate role, conferring the dignity of the "good wife," the "dutiful daughter-in-law," and binding her the more securely by the very esteem it grants. This is the deepest difficulty in the chapter, because it shows that recognition is not simply the cure for the field's wrongs: a recognition that affirms one's place in the hierarchy is not liberation but its gentler enforcement. The criterion that separates affirming recognition from ideological recognition is whether the recognized role remains \emph{contestable}---whether one is recognized as a person who could refuse the role, or recognized only as the perfect occupant of it. The asymmetric field tends to offer the second and call it honor.

\subsection{The distribution of the right to interpret, and the knot}
The third register completes the integration, and it is the one \pinkref{Paper~VI} prepared. Who holds the authority to say what a gift, a silence, an absence \emph{meant}? This is the diplomatic variant of epistemic injustice \citep{fricker2007}, and in the asymmetric field the right of interpretation is distributed along the gradient of power (\xref{sec:power}): the superior party's reading stands as the fact of the matter, the subordinate's as mere subjective impression, to be corrected. The new wife's reading of a slight is "oversensitivity"; the family's reading of her conduct is simply what happened. To suffer one's interpretations systematically discounted is a distinct wrong---hermeneutical, not distributive---a wrong in one's capacity as a knower and a reader of one's own life.

The chapter's integrating claim is that these three are not three separate problems but a single \emb{knot}, and that the hermeneutical strand is what draws it tight. For the party who cannot get their reading of events to count cannot effectively contest the distribution of debt (they cannot establish that the great gift was domination, since their reading of it is discounted), nor the verdict of recognition (they cannot argue that the honor done them is a cage, since to say so is dismissed as ingratitude). Hermeneutical injustice thus \emph{disables the contestation} of the other two: it is the strand that, pulled, tightens the whole. Justice in the diplomatic field is therefore not distributive justice plus recognition plus epistemic fairness, three audits run in parallel; it is one structure in which the maldistribution of the right to interpret locks the maldistribution of debt and the miscarriage of recognition into place. And the knot is tied tightest exactly where the field is most unequal---which is to say that the integration this chapter performs is not a tidy synthesis but a diagnosis of why injustice in intimacy is so resistant to redress: each strand defends the others, and the strand that would let one name the wrong is the first the asymmetry removes.

\goldrule
\footnotesize
\begin{longtable}{Y Y Y Y}
\caption{Value, debt, and justice: literature}\label{tab:justice}\\
\toprule
\rowcolor{pageblush}
Source & Core contribution & Significance for the programme & Limitation / controversy \\
\midrule
\endfirsthead
\toprule
\rowcolor{pageblush}
Source & Core contribution & Significance for the programme & Limitation / controversy \\
\midrule
\endhead
Fricker, \emph{Epistemic Injustice} (2007) & Testimonial / hermeneutical injustice & The distribution of the right to interpret & Scope of the hermeneutical \\
Honneth, \emph{The Struggle for Recognition} (1992) & Recognition vs.\ disrespect & Recognition vs.\ commodification & Ideological recognition \\
Distributive justice (Rawls and after) & Principles of distribution & The distributive face of debt & Abstraction from relation \\
\bottomrule
\end{longtable}

% ============================================================
\section{Conclusion: A Literature Map and a Roadmap}\label{sec:roadmap}

The central claim, restated now that its grounds are in place: diplomacy is not the exception of crisis but the \emb{normal form} in which Generalized Generative Relational Being persists in the world. The dyad does not first exist and then, occasionally, conduct diplomacy; it exists \emph{as} a unit only in continual commerce with the field that surrounds it. To be a "we" in the world is to be, always already, in diplomacy. The paper's arc was the demonstration of this: that the external relational field is constitutive, not additional (\pinkref{§1}); that it must be surveyed without a central phenomenon, by coupling rather than by a hub (\pinkref{§2}); that beneath the field lies a Real the Symbolic cannot exhaust (\xref{sec:real}); that the Symbolic is a tilted field of power (\xref{sec:power}), a grammar of rites and value (\xref{sec:rite}), a hard grammar of law (\xref{sec:law}), and a machinery in which value is made, circulates, and leaks (\xref{sec:polecon}); that a computational-neural bridge carries the field's structure into the diplomatic act (\xref{sec:bridge}); that the act is enacted first in the triad and then across the network (\xref{sec:triad}--\xref{sec:network}); and that justice in this field is a knot of distributive, recognitive, and hermeneutical strands, tied tightest where the field is most unequal (\xref{sec:justice}).

Three threads run the length of that arc and are worth naming as the paper's distinctive contributions, since each is a seed rather than a finished plant. The first is the \emb{re-description of intimacy's relation to the world as diplomacy}, and specifically as \emph{asymmetric} diplomacy---a re-description that lets the tools of value, power, and recognition be brought to bear on a domain usually left to sentiment. The second is the \emb{formal interface}: the proposal that the irreversible costs of an unequal field, and the leakage of value across it, have a precise statement as holonomy, non-conservation around a circuit---a proposal stated minimally here and owed in full to a later paper. The third is the \emb{computational-neural chain} of \xref{sec:bridge}, which threads latent relational state, Bayesian inference, and sequential decision into a single mechanism linking the Symbolic to the act. None of the three is completed here; that is the nature of a prolegomenon. This paper has only surveyed; it has not usurped (in keeping with \pinkref{§2}). What it offers is terrain and coordinates---and, in the tables, an infrastructure---for the work to come, to which the rest of this conclusion points.

\subsection{The total literature map}
The eleven tables of this paper compose, taken together, a single map of the literature a researcher entering this domain must command, clustered by problem-domain and marked, in the difficulty column, by the order in which a newcomer should approach each source. The map is meant to be used in a definite way. A researcher entering one domain---say, the political economy of intimate value---begins from that domain's table, ascends through its sources in order of difficulty (the introductory framing before the specialist treatment), and then, by the "coupling edge" column, locates the adjacent domains that table touches and the tables that treat them, so that the literature is traversed along the couplings rather than as a flat list. The map is thus itself organized by coupling, in keeping with the whole; it is a relational network of readings, not a ranked bibliography. A consolidated synoptic table, merging the eleven into one navigable whole with cross-domain edges drawn explicitly, is the natural appendix to this paper and is left as a discrete deliverable. \todo{assemble the consolidated synoptic table as an appendix}

\subsection{Interfaces to the sub-series to come}
The strata and tables above mark out, with some definiteness, the papers that should follow. Each takes up one node or one coupling edge of the map:
\begin{itemize}
\item A paper on \emb{the power topology of intimate relationships} (developing \xref{sec:power})---the critical keystone, treating the axes of power and their superposition in full.
\item A paper on \emb{the holonomy formalization and the redemption of Value Foam in the diplomatic field} (developing the formal interfaces of \xref{sec:frame}, \xref{sec:power}, \xref{sec:polecon})---the mathematical core, formalizing irreversible deformation and unconserved value.
\item A paper on \emb{the generation of a cross-cultural ``third set of rites''} (developing \xref{sec:thirdrites})---the generative problem, asking whether a shared grammar can be made without reproducing prior asymmetry.
\item A paper on \emb{the psychoanalysis of the in-law triangle} (developing \xref{sec:triad})---transference, the carried real, the mediator's position.
\item A paper on \emb{a computational model of the state inference and decision of relational value} (developing \xref{sec:bridge})---the POMDP of diplomacy, the gift as costly signal, made formal.
\item A paper on \emb{the double grammar of rite and law and their conversion} (developing \xref{sec:law})---which obligations are justiciable, what moves them across the boundary.
\end{itemize}
\todo{fix working titles and ordering of the sub-series}

% ---------- Acknowledgements ----------
\newpage
\section*{\textcolor{rosedeep}{Acknowledgements}}

This paper was written with, and for, one person. To her---of a background in political economy, a lover of travel and of culture, the gentle ``forest girl'' whom the author loves---this work is owed. It was made carrying the thought of her; and it is the author's hope to carry the theoretical framework set out here into the world, to meet, for her sake, every relational difficulty that the real may bring. That hope is the paper's true origin, of which the foregoing theory is only the elaboration.

\vspace{1.2em}
\begin{center}
\begin{minipage}{0.8\linewidth}
\centering
{\cjk\itshape\color{rose}
愿天下有情人，受众目而不为众目所夺，立众缘之中而自成其缘；外睦于亲，内安于己，相视如初。}\\[0.8em]
{\itshape\color{rose}
May lovers in this world bear the eyes of others without being seized by them, and standing amid the world's many bonds, form a bond of their own; at peace with kin without, at peace with themselves within, regarding one another as at the first.}
\end{minipage}
\end{center}

% ---------- References ----------
\newpage
\goldrule
{\color{rose}
\renewcommand{\refname}{\textcolor{rosedeep}{References}}
\small
\bibliographystyle{plainnat}
\bibliography{refs}
}

\end{document}
